\documentclass[12pt,a4paper]{article}

\usepackage[utf8]{inputenc}
\usepackage[T1]{fontenc}
\usepackage{amsmath,amssymb,amsthm}
\usepackage{mathtools}
\usepackage{xcolor}
\usepackage[backref=page]{hyperref}
%\usepackage[backend=biber]{biblatex}
% filecontents is now built into LaTeX kernel


% Theorem environments
\newtheorem{theorem}{Theorem}[section]
\newtheorem{lemma}[theorem]{Lemma}
\newtheorem{proposition}[theorem]{Proposition}
\newtheorem{corollary}[theorem]{Corollary}
\theoremstyle{definition}
\newtheorem{definition}[theorem]{Definition}
\newtheorem{remark}[theorem]{Remark}

% Math operators
\DeclareMathOperator{\Sw}{Sw}
\DeclareMathOperator{\coker}{coker}
\DeclareMathOperator{\Hom}{Hom}
\DeclareMathOperator{\Ext}{Ext}
\DeclareMathOperator{\Ind}{Ind}
\DeclareMathOperator{\Res}{Res}
\DeclareMathOperator{\rk}{rk}
\DeclareMathOperator{\trdeg}{trdeg}
\DeclareMathOperator{\Gal}{Gal}
\DeclareMathOperator{\Aut}{Aut}
\DeclareMathOperator{\GL}{GL}
\DeclareMathOperator{\PGL}{PGL}
\DeclareMathOperator{\id}{id}
\DeclareMathOperator{\Char}{char}
\DeclareMathOperator{\Tor}{Tor}
\DeclareMathOperator{\Br}{Br}
\DeclareMathOperator{\coh}{H}

\title{\bfseries Rationality of the center of a generic division algebra with a
group action over a field of characteristic \(0\)}
\author{Zhengyao Wu\\
  School of Mathematics and Computational Science,\\
  Xiangtan University, Xiangtan, 411105, Hunan, China\\
  \texttt{wuzhengyao@xtu.edu.cn}}
\date{\today}

\begin{document}
\maketitle

\tableofcontents

\begin{abstract}
Let \(F\) be a field of characteristic \(0\) and let
\(H = C_3 \times C_3\) act regularly as a transitive subgroup of \(S_9\).
We resolve Problem~5.2 of Auel--Brussel--Garibaldi--Vishne:
\(Z_H(F,9)\), the center of the generic division algebra of degree~\(9\)
with \(H\)-action, is \emph{not} stably rational and \emph{not} rational,
but is retract rational over~\(F\).  The non-stable-rationality is proved
by computing \(\coh^2(H, M|_H) \cong C_9\) for the restricted Procesi lattice
and showing that its exponent~\(9\) is incompatible with stable permutation, all
permutation lattices have \(\coh^2\)-exponent dividing~\(3\). This integral
cohomological obstruction is of a different nature from the classical unramified
Brauer group.
%(\(\Br_{\mathrm{nr}}(F(H)^H/F) \cong C_3\), Proposition~\ref{prop:B0}\);
% 2026-10-01 第十改：原「which governs the Noether problem but not
% \(Z_H(F,9)\)」从句已删——比较说明归附录 app:brauer。)
The proof generalizes to all odd primes~\(p\):
for \(H = C_p \times C_p \subset S_{p^2}\),
one obtains \(\coh^2(H, M|_H) \cong C_{p^2}\) while permutation lattices
have exponent dividing~\(p\), so \(Z_H(F, p^2)\) is not stably rational.
\end{abstract}

\medskip\noindent
\textbf{Keywords:} generic division algebra, stable rationality,
integral group cohomology, flasque resolution.

\medskip\noindent
\textbf{MSC2020:} 14E08, 20J06, 14M20.

%%%%%%%%%%%%%%%%%%%%%%%%%%%%%%%%%%%%%%%%%%%%%%%%%%%%%%%%%%%%%%%%%%%%%%%%%%%
\section{Introduction}
%%%%%%%%%%%%%%%%%%%%%%%%%%%%%%%%%%%%%%%%%%%%%%%%%%%%%%%%%%%%%%%%%%%%%%%%%%%

Let \(F\) be a field of characteristic \(0\).
For an integer \(n \ge 1\) and a subgroup \(H \subseteq S_n\), the
\emph{generic division algebra with \(H\)-action}, denoted \(UD(F,n,H)\),
has center \(Z_H(F,n)\). Problem~5.2 of
Auel--Brussel--Garibaldi--Vishne \cite[Problem~5.2]{ABGV2011} asks:

\begin{quote}
Determine if \(Z_H(F,9)\) is rational, stably rational, or retract rational
over \(F\), where \(H = C_3 \times C_3 \subset S_9\) is transitive.
\end{quote}

Let \(G\) be a finite group and let \(M\) be a \(\mathbb{Z}[G]\)-lattice,
that is, a finitely generated \(\mathbb{Z}\)-free \(\mathbb{Z}[G]\)-module.
Write \(F(M)\) for the fraction field of the group algebra \(F[M]\),
that is, the function field of the torus with character lattice \(M\),
on which \(G\) acts through the \(\mathbb{Z}[G]\)-module structure of
\(M\); write \(F(M)^G\) for its \(G\)-invariant subfield.
For a subgroup \(H\) of \(G\), write \(M|_H\) for \(M\) regarded as a
\(\mathbb{Z}[H]\)-lattice by restriction of scalars.  Cyclic groups are
denoted \(C_n\), the augmentation ideal is
\(I[G] = \ker(\mathbb{Z}[G] \twoheadrightarrow \mathbb{Z})\), where the
augmentation map sends \(\sum n_g g \mapsto \sum n_g\), and \(\exp(A)\)
denotes the exponent of a finite abelian group \(A\).  Write
\(\coh^i(G, M)\) for the \(i\)-th integral group cohomology of a
\(G\)-module \(M\); for \(\mathbb{Z}[G]\)-lattices this coincides with
Tate cohomology in positive degrees \cite[Def.\ 6.2.4]{Weibel1994}.

A field extension \(K/F\) is \emph{rational} if \(K\) is purely
transcendental over \(F\), and \emph{stably rational} if
\(K(t_1, \dots, t_m)\) is rational over \(F\) for some \(m \ge 0\).

% [2026-10-01 第十改] 「unramified Brauer group 定义」段已注释（用户指令：
% 该段到附录 app:brauer 才用，附录有说明即可）。原文存档：
% For a finitely generated field extension \(K/F\), the \emph{unramified
% Brauer group} \(\Br_{\mathrm{nr}}(K/F)\) is the subgroup of \(\Br(K)\)
% consisting of Brauer classes whose restriction to every discrete valuation
% ring \(F \subset R \subset K\) with fraction field \(K\) is trivial.

The rationality problem for fields of the form \(F(M)^G\) has a long
history.
Endo and Miyata \cite[\S1]{EndoMiyata1973b} developed the theory
of quasi-permutation modules and introduced the notion of quasi-rational
extensions, providing an algebraic foundation for the rationality problem. 
Let \(G\) be a finite group, \(K/k\) a Galois extension with group \(G\),
and \(M\) a \(\mathbb{Z}[G]\)-lattice.  Write \(K(M)\) for the fraction
field of the group algebra \(K[M]\), on which \(G\) acts through the
\(\mathbb{Z}[G]\)-module structure of \(M\) and the Galois action on
\(K\).  They proved \cite[Theorem~1.6]{EndoMiyata1973a} that \(K(M)^G\)
is \emph{quasi-rational} over \(k\) if and only if \(M\) is a
quasi-permutation \(\mathbb{Z}[G]\)-lattice.
Voskresenskii \cite[Chapter~2]{Vos98} gave a comprehensive
treatment via the birational geometry of algebraic tori. 
Colliot-Th\'el\`ene and Sansuc \cite[\S1]{CTS77} developed flasque
resolutions and used them to characterize stable rationality of such
fields in general by the vanishing of the flasque class. 
Saltman \cite[Section~3]{Saltman1984} introduced retract rationality and
constructed the first examples of non-rational stably rational fields. 

When \(H \subset S_n\) is cyclic and transitive,
\(Z_H(F,n)\) is rational if \(F\) contains a primitive \(n\)-th root of unity
(and retract rational for general~\(F\) when \(n\) is not of the form~\(8m\))
\cite[pp.\ 234]{ABGV2011}.
Thus \(H = C_3 \times C_3\) is the first case where the rationality
question is genuinely open for \(n\) odd.

In this paper we give a complete answer:
\begin{table}[tbp]
\centering
\begin{tabular}{lcp{5.5cm}}
\hline
\emph{Property} & \emph{Status} & \emph{Method} \\
\hline
Stably rational & \textbf{No (char~0)}
& \(\coh^2(H, M|_H) \cong C_9\) \\
Rational & \textbf{No (char~0)}
& rational \(\Rightarrow\) stably rational \\
Retract rational & \textbf{Yes (any field)}
& Saltman \cite[Corollary~3.13]{Saltman1984} \\
\hline
\end{tabular}
\caption{Status of the rationality properties of \(Z_H(F,9)\) for
\(H = C_3 \times C_3\).}
\label{tab:main}
\end{table}

The main new ingredient is the computation of the second integral cohomology
group of the restricted Procesi lattice \(M|_H\):
\(\coh^2(H, M|_H) \cong C_9\) (Lemma~\ref{lem:H2-MH}).  The exponent~\(9\) of
this group directly obstructs stable permutation: we prove that for any
permutation \(\mathbb{Z}[H]\)-lattice \(P\) (Definition~\ref{def:lattice}),
every element of \(\coh^2(H, P)\) has order dividing~\(3\)
(Proposition~\ref{prop:perm-exponent}).
Consequently \(M|_H\) cannot be stably permutation
(Definition~\ref{def:lattice}), and by
Theorem~\ref{thm:endo-miyata},
\(Z_H(F,9)\) is not stably rational over \(F\)
(Theorem~\ref{thm:not-stably-rational}).  This integral cohomological
obstruction appears to be new for \emph{multiplicative invariant fields},
i.e.\ the fields \(F(M)^G\) above.

The logical chain is:
\[
\begin{aligned}
&\coh^2(H, M|_H) \cong C_9 \\
&\;\Longrightarrow\; M|_H \text{ is not stably permutation}
  \qquad\text{(Proposition~\ref{prop:stable-exp})} \\
&\;\Longrightarrow\; Z_H(F,9) \text{ is not stably rational}
  \qquad\text{(Theorem~\ref{thm:endo-miyata}, assumed)} \\
&\;\Longrightarrow\; Z_H(F,9) \text{ is not rational}
  \qquad\text{(rationality } \Rightarrow \text{ stable rationality)}.
\end{aligned}
\]
Thus both non-stable-rationality and non-rationality hold over any field
of characteristic~\(0\) (via Theorem~\ref{thm:endo-miyata}, used as a
hypothesis), without reliance on its converse.  The proof generalizes
verbatim to all odd primes~\(p\) (Theorem~\ref{thm:Cp-general}).

% [2026-10-01 第九改] 「As complementary evidence」段已注释（用户指令：
% complementary 说明不重要）。原文存档：
% As complementary evidence, we compute
% \(\Ext^1_{\mathbb{Z}[H]}(I[H], M|_H) \cong C_9 \neq C_3 \times C_3 \cong
% \Ext^1_{\mathbb{Z}[H]}(I[H], 8\cdot\mathbb{Z}[H]\oplus\mathbb{Z})\)
% (Proposition~\ref{prop:M_H-not-permutation}), which provides an independent
% verification that \(M|_H\) is not a permutation lattice.
%
% [2026-10-01 第九改] 「Finally, we prove」段已注释（用户指令：retract
% rationality 系 Saltman 定理所证，表格 tab:main 已提）。原文存档：
% Finally, we prove that \(Z_H(F,9)\) is retract rational over \(F\)
% (Theorem~\ref{thm:retract}), via Saltman's retract rationality theorem
% \cite[Corollary~3.13]{Saltman1984}.

\medskip\noindent
\textbf{Hypotheses.}
The original statement in \cite[Problem~5.2]{ABGV2011} is made over an
arbitrary field.  The cohomological computations
(\(\coh^2(H, M|_H) \cong C_9\) and \(\exp(\coh^2(H, P)) \mid 3\))
are purely integral, valid whenever \(\Char F \nmid 9\).
The geometric conclusions (non-stable-rationality, non-rationality) require
\(\Char F = 0\) (the criterion is stated in characteristic \(0\)):
Theorem~\ref{thm:endo-miyata}, used as a hypothesis.  Retract rationality
holds over any field
(Theorem~\ref{thm:retract}).  Neither non-stable-rationality nor
non-rationality relies on the open converse Endo--Miyata implication.

% [2026-10-01 第九改] 原 \textbf{Generalizability.} 段已前移（用户指令）：
% 去粗体标签，句子并入上方 logic chain 段末（同段）。

% [2026-10-01 第九改] 「Relation to the unramified Brauer group」段已注释
% （用户指令：附录 app:brauer 有）。原文存档：
% \textbf{Relation to the unramified Brauer group.}
% Our obstruction is a purely integral invariant of the lattice \(M|_H\),
% not a geometric invariant of the function field, and the two behave
% differently here: for the regular representation,
% \(\Br_{\mathrm{nr}}(F(H)^H/F) \cong C_3\) and \(F(H)^H\) is not even retract
% rational, whereas \(Z_H(F,9)\) is retract rational but not stably rational,
% so its non-stable-rationality is not detected by \(\Br_{\mathrm{nr}}\).
% Appendix~\ref{app:brauer} develops this comparison, and also places
% \(Z_H(F,9)\), of transcendence degree \(73\), outside the range of the
% existing tables of unramified Brauer groups.

\medskip\noindent
\textbf{Notation.} For the main case (Problem~5.2),
\(H = C_3 \times C_3 = \langle a,b \mid a^3 = b^3 = 1,\; ab = ba \rangle\);
the generalization to \(H = C_p \times C_p\) is treated in
\S\ref{sec:generalization}.
The base field \(F\) is of characteristic \(0\).
For a finite group \(G\), a \(G\)-lattice is a finitely generated
\(\mathbb{Z}[G]\)-module that is \(\mathbb{Z}\)-free. The group ring is denoted
\(\mathbb{Z}[G]\).

\medskip\noindent
\textbf{Ethical and legal declarations}: 

% [2026-10-01 第二十七改，用户指令] Code availability 与 Data availability
% 重复：本节删除（清单句整段并入 Data availability 原位置；链接全篇仅留
% 该处 1 次）。原沿革注释（第二十三／二十六改）见 pre27 逐字节原件。

\medskip\noindent
\textbf{Acknowledgments}. 
The author thanks Xiangtan University for providing the working environment.

\medskip\noindent
\textbf{Research ethics}. Not applicable.

\medskip\noindent
\textbf{Informed consent}. Not applicable. 

\medskip\noindent
\textbf{Author contributions}. The author has accepted responsibility for the
entire content of this manuscript and approved its submission.

\medskip\noindent
\textbf{Use of Large Language Models, AI and Machine Learning Tools}. 
DeepSeek v4 was used, under the author's direction and review, to assist
with literature searching, language polishing, and debugging GAP scripts,
and to produce the two verification programs of Appendix~\ref{app:cocycle}
and the machine-checked Lean development, which is conditional on explicit
named hypotheses (see the Data availability statement).

% [2026-10-01 第二十三改 → 第二十六改] 本块（原自「Code availability」段移入）
% 已按用户追加指令合并为单段：AI 使用披露＋两个验证程序与 Lean 的 GitHub
% 落点＋Lean 条件性指针。原披露句与 Lean 段（第二十六改压缩版）见
% raw/ABGV-5.2-resolution.pre26.tex 与上方第二十六改存档。
% [2026-10-01 第二十七改] 段内 URL 移除（链接全篇仅留 Data availability
% 1 次）：改为指向 Data availability 语句。
% [2026-10-01 第二十八改] 按用户指令再精简（两句）。

% [2026-10-01 第七改，用户指令] 新增一段：几何等同归约层（round (1)）的假设表。
% 义务出处：wiki/proj-ABGV52-EM-criterion.md §3an 末节「假设表更新义务（承重）」——
% G 档（= 几何等同归约层）写入后必须补 Procesi 1967，否则读者会误读为
% 「几何等同无条件证出」。原缺载事实：开发中 Z_H(F,p²) ≅ F(M|_H)^H
% 是**具名假设**（ProcesiHyp，见 scripts/lean_ws/ABGV52/GThm_ABGV52.lean），
% 该层已证的是「归约」：正则限制 U|_H ≅ Z[H]、Procesi 复形、π|_H 的识别、
% 以及不变量域沿该识别的搬运（G1c/G1d/G2p）。
% [2026-10-01 第二十六改，用户指令] 「Lean 形式化说明」再压缩为一段（数行），
% 且与 repo README 分工不重复：论文段只留存在＋条件性＋指针；细节（no sorry／
% 公理名／假设表／文件表／构建说明）归 README，论文不复述。
% （追加指令 2026-10-01：gap／python 代码说明并入 LLM 段，全节合并为单段。）
% 被替换块全文（第二十五改压缩版）存档如下；更早原始句另见上方
% 第二十五改存档 (a)–(e) 与第七改注释。逐字节原件另存
% raw/ABGV-5.2-resolution.pre26.tex。
% --- 原第二段＋第三段（状态／假设／无条件部分／Cp-general）---
% \medskip\noindent The Lean development verifies the logical chain of the main
% theorem. It compiles with no \texttt{sorry} and depends only on the three
% standard axioms \texttt{propext}, \texttt{Classical.choice} and
% \texttt{Quot.sound}. One input is taken as an \emph{explicit hypothesis}: the
% Endo--Miyata--Voskresenskii criterion (Theorem~\ref{thm:endo-miyata}), which
% the main text also assumes. Everything else, including the computation
% \(\coh^2(H, M|_H) \cong C_9\) of Lemma~\ref{lem:H2-MH}, the implication of
% Proposition~\ref{prop:stable-exp}, and the faithfulness of the action in
% Proposition~\ref{prop:faithful-MH}, is proved inside the development rather
% than assumed, so that Corollary~\ref{thm:not-stably-rational} is formalized
% with the criterion as its only hypothesis.
%
% \medskip\noindent Theorem~\ref{thm:Cp-general} is formalized with the odd
% prime \(p\) as a parameter, the case \(p = 3\) used above being the
% corresponding instance; its stable-rationality conclusion again carries the
% criterion as its only hypothesis, and statement~(4) is reduced to
% \cite[Corollary~3.13]{Saltman1984}, taken as an explicit hypothesis rather
% than formalized (the same reduction underlies Corollary~\ref{thm:retract}).
% --- 原几何段（含条件性声明；第七改注释所述假设表段）---
% \medskip\noindent The geometric identification \(Z_H(F, p^2) \cong F(M|_H)^H\)
% of Remark~\ref{rem:procesi-geometric} (Procesi \cite{Procesi1967}) is used in
% the development as a \emph{named hypothesis}: the formalization of that layer
% establishes the reduction to the regular restriction
% (Proposition~\ref{prop:U-restrict-canonical}), the Procesi exact sequence of
% Proposition~\ref{prop:procesi-H}, and the transport of \(H\)-invariants along
% the identification, but neither constructs the field \(Z_H(F, p^2)\) nor
% proves Procesi's theorem. The development is therefore a \emph{conditional}
% formalization of that logical chain, not a formalization of the theorems. The
% field \(Z_H(F,9)\) itself is not constructed in the development; the corollary
% is formalized for an abstract field linked to \(M|_H\) by the criterion, and
% the cohomological computation is certified numerically in
% Appendix~\ref{app:cocycle}. The full development, with a file-by-file
% description and instructions for reproducing the verification, is available at
% \url{https://github.com/wuzhengyao/abgv52-lean}.

% [2026-10-01 第二十五改，用户指令] 长说明压缩为链接说明（上方 repo URL 句）；
% 具名假设（EM 判据／Procesi 几何识别／Saltman 3.13）与验证状态句保留。
% 被压缩的原句存档如下（注释而非删除；文字逐字，行宽重排）：
% 逐字节原件（编辑前全文件）另存 raw/ABGV-5.2-resolution.pre25.tex。
% （a）The computation \(\coh^2(H,M|_H) \cong C_9\) (Lemma~\ref{lem:H2-MH})
% is \emph{not} an input: its proof is reproduced inside the development,
% together with the discriminant pair it yields.
% （b）The implication of Proposition~\ref{prop:stable-exp} is also proved
% inside the development, not assumed: it is derived by combining Shapiro's
% lemma, applied to the orbit decomposition of a finite \(H\)-set, with the
% short exact sequence \(0 \to \mathbb{Z} \to \mathbb{Q} \to
% \mathbb{Q}/\mathbb{Z} \to 0\) of trivial modules over the corresponding
% subgroup \(J \le H\).
% （c）The faithfulness of the action of \(H\) on \(M|_H\)
% (Proposition~\ref{prop:faithful-MH}) is also proved inside the
% development, by the coordinate argument of the second proof of that
% proposition, which does not use the Procesi identification.
% （d）Theorem~\ref{thm:Cp-general} is formalized as well, with the odd
% prime \(p\) as a parameter: for every odd prime \(p\), the development
% proves that \(\coh^2(H, M|_H) \cong C_{p^2}\) and that \(\coh^2\) of a
% permutation lattice has exponent dividing \(p\), and hence that \(M|_H\)
% is not a stably permutation lattice; the case \(p = 3\) used above is
% the corresponding instance.
% （e）Consequently Corollary~\ref{thm:not-stably-rational} is formalized
% with the criterion as its only hypothesis.
% [2026-09-30 移除，原句存档] The development also assumes, as an additional
% hypothesis, the implication of Proposition~\ref{prop:stable-exp}, whose
% proof in the main text is not formalized; the remaining steps are proved.
% 移除理由：prop:stable-exp 的蕴含已由 Lean 开发证明（scripts/lean_ws/ABGV52/
% S4_ABGV52.lean 的 stable_exp；S5_ABGV52.lean 的去假设接线版
% not_stably_rational_wired' 假设表只剩 EM 判据一条）。
% [2026-09-30 补记] 增补两句：① 点名 \ref{thm:not-stably-rational} 的假设表
% （判据一条）；② 负向范围声明——Z_H(F,9) 在开发中未构造，该推论是对
% 「抽象域 ＋ 判据接口」陈述的。原句「With the EM criterion as its only
% remaining hypothesis, ...」改作「The development is therefore ...」，
% 以免与 ① 重复（信息未删，仍在段内）。
% [2026-09-30 第五改] 新增一段：\ref{thm:Cp-general} 的形式化（奇素数 p 为参数）——
% 无条件部分为 coh²(H,M|_H) ≅ C_{p²} 与置换格的指数界（⟹ M|_H 非稳定置换）；
% 稳定有理性侧仍以 EM 判据为唯一假设；第 (4) 条归约为 \cite[Cor 3.13]{Saltman1984}
% 并作显式假设。p=3 特例按用户裁定**不单列**（一般陈述的实例）。
% [2026-09-30 第六改，用户选「乙」] 同段「范畴记号 ＋ 字母撞车」两病并治。
% 原句存档（下三行）：
%   "... it is derived from the short exact sequence
%    0 -> Z -> Q -> Q/Z -> 0 in Rep_Z K, together with Shapiro's lemma
%    applied to the orbit decomposition of a finite H-set."
% 两病：① `Rep_Z K` 是 **Lean 侧范畴记号**，正文从未定义（全篇仅此一处）；
% ② 该处 `K` 是**群**（子群），而正文通篇 `K` 是**域**（K/F、K/k、L/K）
% ⟹ 同字母换义。处置：去记号（改叙述 "trivial modules over the
% corresponding subgroup"）＋ 该子群**改名 J ≤ H**，与域 K 彻底脱钩。

\medskip\noindent
\textbf{Conflict of interest}. 
The author declares that there is no conflict of interest. 

\medskip\noindent
\textbf{Research funding}. 
The author is supported by Xiangtan University project
09KZ-KZ08101 \textit{algebraic invariants and local-global principles}.

\medskip\noindent
\textbf{Data availability}. 
The two verification programs of Appendix~\ref{app:cocycle}, their
recorded output for \(p=3\) and \(p=5\), a \texttt{SHA-256} manifest,
and the machine-checked Lean development are collected in the
accompanying code package, available at
\url{https://github.com/wuzhengyao/abgv52-lean}.

%%%%%%%%%%%%%%%%%%%%%%%%%%%%%%%%%%%%%%%%%%%%%%%%%%%%%%%%%%%%%%%%%%%%%%%%%%%
\section{Preliminaries}
%%%%%%%%%%%%%%%%%%%%%%%%%%%%%%%%%%%%%%%%%%%%%%%%%%%%%%%%%%%%%%%%%%%%%%%%%%%

\begin{definition}\label{def:lattice}
A \emph{\(\mathbb{Z}[G]\)-lattice} is a finitely generated
\(\mathbb{Z}[G]\)-module that is \(\mathbb{Z}\)-free.
A \(\mathbb{Z}[G]\)-lattice \(M\) is called:
\begin{enumerate}
  \item \emph{permutation} if \(M \cong \bigoplus_{i=1}^r \mathbb{Z}[G/K_i]\)
    for subgroups \(K_i \le G\), where \(\mathbb{Z}[G/K]\) denotes the
    \(\mathbb{Z}[G]\)-module with \(\mathbb{Z}\)-basis the left cosets
    \(G/K\) and \(G\)-action by left multiplication;
  \item \emph{stably permutation} if there exist permutation lattices \(P, Q\)
    such that \(M \oplus P \cong Q\);
% \item \emph{stably isomorphic} to \(N\) (written \(M \sim_{\text{st}} N\)) if
% \(M \oplus \mathbb{Z}[G]^{\oplus a} \cong N \oplus \mathbb{Z}[G]^{\oplus b}\)
% for some \(a,b \ge 0\).
\end{enumerate}
These concepts are standard; see \cite[\S1]{CTS77}.
\end{definition}

\begin{lemma}\label{lem:perm-char}
Let \(H\) be a finite group and \(K \trianglelefteq H\) a normal
subgroup, and let \(\mathbb{Z}[H/K]\) be the permutation
\(\mathbb{Z}[H]\)-lattice
(Definition~\ref{def:lattice}).  The character of the rational
representation \(\mathbb{Q}[H/K]\) is
\[
\chi_{H/K}(g) = \bigl|\{xK \in H/K : g x K = x K\}\bigr|
= \begin{cases}
[H:K], & g \in K,\\
0, & g \notin K.
\end{cases}
\]
\end{lemma}

% [2026-10-01 第二十改，用户指令] 原第二半句（含 CR Ex 10.4 引）已删——
% 「交换 H 时重数 m_K 由特征标唯一决定」为**假**：机器核反例（KB 侧
% `PC2_ABGV52.lean` 的 `permChar_uniqueness_false`：3 单点⊔正则 与 四个 C₃-型
% 轨道两 H-集，字符逐元素相同而平凡轨道重数 3 ≠ 0、二者不同构）；且所引
% Curtis-Reiner Example 10.4 原文只是「诱导模 ≅ 置换模」的识别，不含唯一性。
% 原文存档：
% Moreover, for an abelian group~\(H\), the multiplicities
% \(m_K\) in a permutation lattice
% \(P \cong \bigoplus_{K \le H} m_K \cdot \mathbb{Z}[H/K]\)
% are uniquely determined by the character of
% \(P \otimes_{\mathbb{Z}} \mathbb{Q}\)
% \cite[Example~10.4]{CurtisReiner1981}.

This lemma is the key tool in Lemma~\ref{lem:perm-class} below, where
the rational character of a permutation lattice yields linear
equations for its multiplicities.

% [2026-10-01 第二十一改，用户指令] flasque 整块（def:flasque ＋
% def:flasque-class-group ＋ prop:flasque-resolution）移至附录 C
% （「Unramified Brauer group」段之前；该块全篇唯一实质使用者——
% 两个定义仅被该命题引用，命题仅被附录 C 引用）。原文存档如下：
%
% \begin{definition}\label{def:flasque}
% A \(\mathbb{Z}[G]\)-lattice \(F\) is \emph{flasque} if
% \(\hat{\coh}^{-1}(K, F) = 0\) for all subgroups \(K \le G\)
% (equivalently, \(\Ext^1_G(F, P) = 0\) for every permutation lattice~\(P\))
% \cite[\S1, Lemme~9]{CTS77}. Here \(\hat{\coh}^{-1}\) denotes
% Tate cohomology in degree \(-1\); for a finite group~\(G\) and a
% \(\mathbb{Z}[G]\)-lattice~\(M\),
% \(\hat{\coh}^{-1}(G, M) = \ker(N_G) / I[G]M\), where
% \(N_G = \sum_{g \in G} g\) is the norm element and
% \(I[G] = \ker(\mathbb{Z}[G] \twoheadrightarrow \mathbb{Z})\) is the
% augmentation ideal.
% \end{definition}
%
% \begin{definition}\label{def:flasque-class-group}
% Two \(\mathbb{Z}[G]\)-lattices \(M, N\) are called \emph{similar}
% (written \(M \sim N\)) if \(M \oplus S_1 \cong N \oplus S_2\) for
% some permutation lattices \(S_1, S_2\) (Definition~\ref{def:lattice}).
% This is an equivalence relation; denote the similarity class of \(M\)
% by \([M]\). Let
% \[
% C(G)/S(G) = \{ [M] : M \text{ a finitely generated }
% \mathbb{Z}[G]\text{-lattice} \},
% \]
% a commutative monoid under direct sum, with zero element the class
% of permutation lattices. Its maximal subgroup is
% \[
% D(G)/S(G) = \{ [M] \in C(G)/S(G) : M \text{ is stably permutation}\},
% \]
% which is the semigroup \(C(G)/S(G)\) of \cite[\S4.7]{Vos98}.
% \end{definition}
%
% \begin{proposition}\label{prop:flasque-resolution}
% Every \(\mathbb{Z}[G]\)-lattice~\(M\) has a \emph{flasque resolution}
% \[
% 0 \longrightarrow M \longrightarrow P \longrightarrow F \longrightarrow 0
% \]
% with \(P\) permutation (Definition~\ref{def:lattice}) and \(F\) flasque
% (Definition~\ref{def:flasque})
% \cite[\S1, Lemme~3]{CTS77}. The similarity class
% \([F] \in D(G)/S(G)\) (Definition~\ref{def:flasque-class-group})
% depends only on \(M\)
% \cite[\S1, Lemme~5]{CTS77}; we denote it by
% \(\rho(M) = [F]\) and call it the
% \emph{flasque class} of \(M\).
% By construction, \(M\) is stably permutation iff \(\rho(M) = 0\)
% \cite[\S1, Lemmes~7 and~8]{CTS77}; this is the defining property
% of the group \(D(G)/S(G)\).
% \end{proposition}

\begin{theorem}\label{thm:endo-miyata}
Let \(G\) be a finite group and \(M\) a \(\mathbb{Z}[G]\)-lattice on
which \(G\) acts faithfully.  Let \(F\) be a field of characteristic
\(0\).  Here \(F(M)\) and \(F(M)^G\) are as in the Introduction, and
\(G\) acts on \(F(M)\) through the \(\mathbb{Z}[G]\)-module structure
of~\(M\).  If \(F(M)^G\) is \emph{stably rational} over \(F\), then
\(M\) is a stably permutation \(\mathbb{Z}[G]\)-lattice.  In the
literature this criterion is stated for the \emph{twisted}
invariant field \(Q(L/K,M) = L(M)^G\)
of \cite[Theorem~1.3(b)]{Saltman1987}, where \(L/K\) is Galois with
group \(G\) and \(G\) acts on the coefficient field as well; the
field \(F(M)^G\) above is the field \(F(M,G)\) of
\cite[p.~221]{Saltman1987}, and for it the criterion is an
assumption rather than a quotation of a proved result.
\textbf{We use this criterion as a hypothesis}; it is not proved in
this paper.
% [2026-09-30 引文勘误] 原句尾挂 \cite[\S1, Lemme~6]{CTS77}，已移除。
% 核据：CTS77 印刷 Lemme 6（印刷页 181）逐字为「Les applications ρ et ζ
% sont additives et duales l'une de l'autre … ρ est un épimorphisme de S_G
% sur F_G」，不含任何域/有理性陈述；该判据在文献中无出处，与本定理随后
% 的自陈（an assumption rather than a quotation）自相矛盾。可用邻近锚 =
% §1 Lemme 8（印刷页 182，ρ 相等的判据）与 §2 Proposition 6（印刷页 189，
% 环面情形：ρ(T) 对 k-稳定有理的 T 平凡，且 ρ(T) = ρ(hat T)）。
% [2026-10-01 第二十一改，用户指令] 原句尾「, i.e., its flasque class
% \(\rho(M) \in D(G)/S(G)\) vanishes」已删——flasque 类术语连同其定义
% （def:flasque／def:flasque-class-group）移入附录 C；判据逻辑不变（措辞清理）。
\end{theorem}

% [2026-10-01 第十八改，用户指令] `thm:endo-miyata-converse` 主线不使用
% （各处提及均仅为「不依赖它」的声明；其引用仅 Vos98 一处），注释存档如下：
%
% \begin{theorem}\label{thm:endo-miyata-converse}
% Let \(G\) be a finite \emph{abelian} group and \(M\) a
% \(\mathbb{Z}[G]\)-lattice on which \(G\) acts faithfully.  Let \(F\) be
% a field of characteristic \(0\) containing a primitive
% \(\exp(G)\)-th root of unity.  Then \(F(M)^G\) is \emph{stably rational}
% over \(F\) if \(M\) is a stably permutation \(\mathbb{Z}[G]\)-lattice
% \cite[\S7.1, (1) and \S7.2]{Vos98}.
% \end{theorem}

In the sequel we apply Theorem~\ref{thm:endo-miyata} in the
contrapositive direction, the criterion being used as a hypothesis
throughout: if \(M|_H\) is \emph{not} stably permutation, then
\(F(M|_H)^H \cong Z_H(F,9)\) is \emph{not} stably rational.
The cohomological computation of \S\ref{sec:cohom} proves that
\(\coh^2(H, M|_H) \cong C_9\) while permutation lattices have
\(\coh^2\)-exponent dividing~\(3\); hence \(M|_H\) is not stably
permutation (Theorem~\ref{thm:not-stably-perm}), and the criterion
yields non-stable-rationality (Corollary~\ref{thm:not-stably-rational}).
The same strategy generalizes to \(C_p \times C_p\) in
\S\ref{sec:generalization}.

% [2026-10-01 第十七改，用户指令] `thm:UCT` 移至附录 B（原处仅此一处定义，
% 其使用全部在附录 B/C：`prop:B0` 证明与附录 C 泛化节）；原文存档如下：
%
% \begin{theorem}\label{thm:UCT}
% For a finite group \(G\) and a trivial \(G\)-module \(M\), there is a
% natural short exact sequence
% \[
% 0 \longrightarrow \Ext^1_{\mathbb{Z}}(\mathrm{H}_{n-1}(G, \mathbb{Z}), M)
% \longrightarrow \coh^n(G, M)
% \longrightarrow \Hom(\mathrm{H}_n(G, \mathbb{Z}), M) \longrightarrow 0
% \]
% for all \(n \ge 1\) \cite[3.6.5]{Weibel1994}.
% If \(M\) is divisible (e.g.\ \(M = \mathbb{Q}/\mathbb{Z}\)), the
% sequence splits, giving a direct sum decomposition.
% \end{theorem}
%
% This theorem is used in Proposition~\ref{prop:B0} to compute
% \(\coh^2(H, \mathbb{Q}/\mathbb{Z})\).

% \begin{theorem}[Homology of abelian groups]\label{thm:H2-abelian}
% Let \(k\) be a principal ideal domain and \(G\) an abelian group.
% There is a natural map
% \[
% \psi \colon \bigwedge\!^*\, (G \otimes k) \longrightarrow H_*(G, k)
% \]
% with the following properties:
% \begin{enumerate}
% \item[(i)] \(\psi\) is injective, and split injective if \(G\) is finitely
% generated.
% \item[(ii)] If every prime \(p\) such that \(G\) has \(p\)-torsion is
% invertible in \(k\), then \(\psi\) is an isomorphism.
% \item[(iii)] If \(k\) has characteristic zero (e.g.\ \(k = \mathbb{Z}\)),
% then \(\psi\) is an isomorphism in dimension~\(2\).
% \end{enumerate}
% \cite[Ch.\ V, Theorem~6.4]{Brown1982}.
% \end{theorem}

\section{The Procesi Lattice}
\label{sec:procesi}
%%%%%%%%%%%%%%%%%%%%%%%%%%%%%%%%%%%%%%%%%%%%%%%%%%%%%%%%%%%%%%%%%%%%%%%%%%%

We give a self-contained construction of the Procesi
\(\mathbb{Z}[S_9]\)-lattice~\(M\) and the Procesi isomorphism, keeping only what
is needed for the rationality computations. An accessible modern exposition is
in \cite[\S5]{ABGV2011} and
\cite{Formanek2002}; the original reference is \cite{Procesi1967}.

\subsection{Construction over \(S_9\)}

Let \(n = 9\) and let \(\Omega = S_9/S_8\) be the set of left cosets of \(S_8\)
in \(S_9\). Then \(|\Omega| = 9\) and \(S_9\) acts on \(\Omega\) by left
multiplication. Define the permutation \(\mathbb{Z}[S_9]\)-module
\[
U = \mathbb{Z}[\Omega] \cong \bigoplus_{\omega \in \Omega} \mathbb{Z}\cdot
e_\omega ,
\]
where \(\{e_\omega\}_{\omega \in \Omega}\) is the standard basis and
\(g \cdot e_\omega = e_{g\omega}\) for \(g \in S_9\).  Consider the
\(\mathbb{Z}[S_9]\)-module map
\[
\varphi : U \otimes_{\mathbb{Z}} U \longrightarrow \mathbb{Z}[\Omega], \qquad
\varphi(e_\alpha \otimes e_\beta) = e_\alpha - e_\beta ,
\]
extended \(\mathbb{Z}\)-linearly.  The map \(\varphi\) is \(S_9\)-equivariant:
\[
\varphi\bigl(g \cdot (e_\alpha \otimes e_\beta)\bigr)
= \varphi(e_{g\alpha} \otimes e_{g\beta})
= e_{g\alpha} - e_{g\beta}
= g \cdot (e_\alpha - e_\beta)
= g \cdot \varphi(e_\alpha \otimes e_\beta).
\]

The image of \(\varphi\) is the \emph{augmentation submodule}
\[
I[\Omega] = \ker\!\Bigl( \varepsilon : \mathbb{Z}[\Omega] \twoheadrightarrow
\mathbb{Z},\;
\varepsilon\bigl(\sum a_\omega e_\omega\bigr) = \sum a_\omega \Bigr).
\]
Indeed, every difference \(e_\alpha - e_\beta\) has augmentation zero, and
conversely any element \(\sum a_\omega e_\omega\) with \(\sum a_\omega = 0\) is
a \(\mathbb{Z}\)-linear combination of such differences (take \(\{e_\alpha -
e_{\omega_0} : \alpha \neq \omega_0\}\) for a fixed \(\omega_0\)). Moreover,
\(\varphi\) is surjective onto \(I[\Omega]\).

\begin{definition}
The \emph{Procesi \(\mathbb{Z}[S_9]\)-lattice} is \(M = \ker(\varphi)\).
\end{definition}

By construction we have the short exact sequence of \(\mathbb{Z}[S_9]\)-modules
\begin{equation}\label{eq:procesi-S9}
0 \longrightarrow M \longrightarrow U \otimes_{\mathbb{Z}} U \xrightarrow{\
\varphi\ } I[\Omega] \longrightarrow 0 .
\end{equation}

The module \(M\) is \(\mathbb{Z}\)-free: it is a \(\mathbb{Z}\)-submodule of the
free \(\mathbb{Z}\)-module \(U \otimes_{\mathbb{Z}} U\), and every submodule of
a free \(\mathbb{Z}\)-module is free (since \(\mathbb{Z}\) is a principal ideal
domain)
\cite[Ch.VII, \S3, Th.1, Cor.2]{Bourbaki2007b}.
The ranks are
\[
\rk_{\mathbb{Z}} U = 9,\qquad
\rk_{\mathbb{Z}}(U \otimes U) = 81,\qquad
\rk_{\mathbb{Z}} I[\Omega] = 8,\qquad
\rk_{\mathbb{Z}} M = 73.
\]

\begin{remark}\label{rem:procesi-geometric}
Procesi \cite{Procesi1967} proved that the center of the generic matrix
algebra satisfies \(Z_{S_9}(F,9) \cong F(M)^{S_9}\), where \(F(M)^{S_9}\) is
the invariant field of the purely transcendental extension \(F(M)\) under
the natural \(S_9\)-action induced by the \(\mathbb{Z}[S_9]\)-module
structure of \(M\).  For any subgroup
\(H \subseteq S_9\), taking \(H\)-invariants gives
\[
Z_H(F,9) = F(M)^H \cong F(M|_H)^H,
\]
This isomorphism holds over any field \(F\) of characteristic~\(0\);
see \cite[\S5]{ABGV2011}, \cite{Formanek2002}, and the references therein.
For the computations that follow, only the exact
sequence~\eqref{eq:procesi-S9} and this identification are required.
\end{remark}

\subsection{Restriction to \(H = C_3 \times C_3\)}

Now fix a transitive embedding \(H = C_3 \times C_3 \hookrightarrow S_9\),
i.e., an embedding whose image acts transitively on
\(\Omega = S_9/S_8\).
Since \(|H| = 9 = |\Omega|\), the action is moreover regular
(i.e., free and transitive, called \emph{simply transitive}
in \cite[ch.I, \S5, no.6, def.7]{Bourbaki2007a}):
choosing a basepoint \(\omega_0 \in \Omega\), the map
\(H \to \Omega\), \(h \mapsto h \cdot \omega_0\), is an \(H\)-equivariant
bijection. Consequently we obtain an explicit isomorphism of
\(\mathbb{Z}[H]\)-modules
\begin{equation}\label{eq:U-restrict}
U|_H \xrightarrow{\ \cong\ } \mathbb{Z}[H], \qquad
e_{h \cdot \omega_0} \longmapsto h \quad (h \in H).
\end{equation}

\begin{proposition}\label{prop:U-restrict-canonical}
The isomorphism~\eqref{eq:U-restrict} is independent of the choice of
\(\omega_0\) up to a \(\mathbb{Z}[H]\)-module automorphism of
\(\mathbb{Z}[H]\).  Consequently,
\(\Ext^1_{\mathbb{Z}[H]}(-, M|_H)\), \(\coh^2(H, M|_H)\), and the
isomorphism class of \(M|_H\) are independent of this choice.
\end{proposition}

\begin{proof}
Let \(\omega_0' = g \cdot \omega_0\) be another basepoint.
The isomorphism induced by \(\omega_0'\) sends
\(e_{h \cdot \omega_0'} = e_{hg \cdot \omega_0} \longmapsto h\),
while the original sends \(e_{hg \cdot \omega_0} \longmapsto hg\);
hence the two isomorphisms differ by right multiplication by
\(g^{-1}\) on \(\mathbb{Z}[H]\).  Right multiplication by a fixed
element of \(H\) is a \(\mathbb{Z}[H]\)-module automorphism of the
left regular module.  Therefore the induced map on cohomology and
\(\Ext^1\) is the identity, and the isomorphism class of \(M|_H\) is unchanged.
\end{proof}

\begin{proposition}\label{prop:embedding-independence}
Let \(\iota_1, \iota_2 : H \hookrightarrow S_9\) be two transitive embeddings.
Then \(M|_{\iota_1(H)} \cong M|_{\iota_2(H)}\) as \(\mathbb{Z}[H]\)-modules.
\end{proposition}

\begin{proof}
Both \(\iota_1\) and \(\iota_2\) make \(\Omega\) into a regular
\(H\)-set (since \(|H| = 9 = |\Omega|\) and each action is transitive).
Any two regular \(H\)-sets of the same cardinality are
\(H\)-equivariantly isomorphic: fix basepoints
\(\omega_1, \omega_2 \in \Omega\), and define a bijection
\(\Omega \to \Omega\) by \(h \cdot \omega_1 \longmapsto h \cdot \omega_2\).
This bijection is an element \(\sigma \in S_9\), and by construction
\(\sigma \, \iota_1(h) \, \sigma^{-1} = \iota_2(h)\) for all \(h \in H\).
Conjugating the embedding by \(\sigma\) replaces the
\(H\)-action on \(U \otimes U\) with its conjugate under the \(S_9\)-action,
which induces an isomorphism \(M|_{\iota_1(H)} \cong M|_{\iota_2(H)}\) of
\(\mathbb{Z}[H]\)-modules via the identification of \(H\) with
\(\iota_1(H)\) and \(\iota_2(H)\) followed by the conjugation
\(\iota_2^{-1} \circ \iota_1\).
\end{proof}

\begin{proposition}\label{prop:tensor-reg-ZH}
There is an isomorphism of \(\mathbb{Z}[H]\)-modules
\[
\mathbb{Z}[H] \otimes_{\mathbb{Z}} \mathbb{Z}[H] \;\cong\; \mathbb{Z}[H]^{\oplus
9},
\]
where the left side carries the diagonal action
\(h \cdot (x \otimes y) = hx \otimes hy\) and the right side
carries the action \(h \cdot (z_b)_{b \in H} = (h z_{h^{-1} b'})_{b' \in H}\),
that is, left multiplication composed with the coordinate permutation
\(b \mapsto hb\).
\end{proposition}

\begin{proof}
Identify \(\mathbb{Z}[H]\) with the induced module
\(\Ind_1^H(\mathbb{Z}) = \mathbb{Z}[H] \otimes_{\mathbb{Z}} \mathbb{Z}\),
where \(\Ind_1^H\) denotes induction from the trivial subgroup~\(\{1\}\)
to~\(H\), and \(\Res_1^H\) denotes restriction from~\(H\) to~\(\{1\}\).
For \(\mathbb{Z}\)-modules \(A, B\) with trivial \(H\)-action,
the projection formula for induced modules
\cite[Corollary~10.20]{CurtisReiner1981} states
\[
\Ind_1^H(A) \otimes_{\mathbb{Z}} \Ind_1^H(B)
\;\cong\; \Ind_1^H\!\bigl( A \otimes_{\mathbb{Z}} \Res_1^H \Ind_1^H(B) \bigr)
\]
as \(\mathbb{Z}[H]\)-modules (diagonal action on the left, induced action
on the right).  Taking \(A = B = \mathbb{Z}\) gives
\[
\mathbb{Z}[H] \otimes_{\mathbb{Z}} \mathbb{Z}[H]
\;\cong\; \Ind_1^H\!\bigl( \Res_1^H \Ind_1^H(\mathbb{Z}) \bigr).
\]
The restriction \(\Res_1^H \Ind_1^H(\mathbb{Z})\) is the regular
module \(\mathbb{Z}[H]\) viewed merely as a \(\mathbb{Z}\)-module
(forgetting the \(H\)-action); by definition \(\mathbb{Z}[H]\) has
\(\mathbb{Z}\)-basis \(\{h : h \in H\}\), so
\(\Res_1^H \Ind_1^H(\mathbb{Z}) \cong \mathbb{Z}^{\oplus 9}\).
Applying \(\Ind_1^H\) then yields
\(\Ind_1^H(\mathbb{Z}^{\oplus 9})
\cong \Ind_1^H(\mathbb{Z})^{\oplus 9}
\cong \mathbb{Z}[H]^{\oplus 9}\).
The isomorphism is \(H\)-equivariant: choose the \(\mathbb{Z}\)-basis
\(\{g\}_{g\in H}\) of the second tensor factor \(\mathbb{Z}[H]\);
the map \(x \otimes g \longmapsto x\) in the \(g\)-th coordinate
satisfies \(h \cdot (x \otimes g) = hx \otimes hg \longmapsto hx\)
in the \(hg\)-th coordinate.
\end{proof}

Under the identification~\eqref{eq:U-restrict} and
Proposition~\ref{prop:tensor-reg-ZH},
the restriction of the map \(\varphi\) to \(H\) yields the following.

\begin{proposition}\label{prop:procesi-H}
There is an exact sequence of \(\mathbb{Z}[H]\)-modules
\begin{equation}
0 \longrightarrow M|_H \longrightarrow \mathbb{Z}[H]^{\oplus 9}
\xrightarrow{\ \pi|_H\ } I[H] \longrightarrow 0 ,
\end{equation}
where \(I[H] = \ker(\varepsilon : \mathbb{Z}[H] \twoheadrightarrow
\mathbb{Z})\), \(\varepsilon(\sum a_h h) = \sum a_h\), is the augmentation ideal
of \(H\). All modules in the sequence are finitely generated and
\(\mathbb{Z}\)-free, and \(\rk_{\mathbb{Z}} M|_H = 73\).
\end{proposition}

\begin{proof}
\emph{Exactness.}
Under the identifications, the map \(\varphi\) becomes a
\(\mathbb{Z}[H]\)-module map
\(\pi|_H : \mathbb{Z}[H]^{\oplus 9} \longrightarrow I[H]\).
Each standard basis element of \(\mathbb{Z}[H]^{\oplus 9}\) maps to
a difference \(h - h'\) in \(\mathbb{Z}[H]\); these differences
generate the \(\mathbb{Z}\)-module \(I[H]\), and \(I[H]\) is generated
as a \(\mathbb{Z}[H]\)-module by any single element \(h-1\) with
\(h \neq 1\); hence \(\pi|_H\) is surjective.
Define \(M|_H = \ker(\pi|_H)\), the restriction of the Procesi
lattice to \(H\).

\medskip\noindent
\emph{Freeness.}
All modules in the sequence are finitely generated and
\(\mathbb{Z}\)-free: \(\mathbb{Z}[H]^{\oplus 9}\) is \(\mathbb{Z}\)-free
with basis the group elements in each coordinate; \(I[H]\) is a
\(\mathbb{Z}\)-submodule of \(\mathbb{Z}[H]\), hence free
\cite[Ch.VII, \S3, Th.1, Cor.2]{Bourbaki2007b};
\(M|_H = \ker(\pi|_H)\) is a \(\mathbb{Z}\)-submodule of
\(\mathbb{Z}[H]^{\oplus 9}\), hence free.

\medskip\noindent
\emph{Rank.}
\(\rk_{\mathbb{Z}} M|_H
= \rk_{\mathbb{Z}} \mathbb{Z}[H]^{\oplus 9}
- \rk_{\mathbb{Z}} I[H]
= 9 \cdot 9 - 8 = 73\),
matching the rank of \(M\) over \(S_9\).
\end{proof}

\begin{proposition}\label{prop:MH-char}
\[
M|_H \otimes_{\mathbb{Z}} \mathbb{Q} \;\cong\; 8\cdot \mathbb{Q}[H] \;\oplus\;
\mathbb{Q},
\]
where \(\mathbb{Q}\) denotes the one-dimensional trivial
\(\mathbb{Q}[H]\)-module.
\end{proposition}

\begin{proof}
Tensor the exact sequence of Proposition~\ref{prop:procesi-H}
with \(\mathbb{Q}\) (flat over \(\mathbb{Z}\),
\cite[Ch.II, \S 2, no.4, Th.1]{Bourbaki2006}):
\[
0 \longrightarrow M|_H \otimes_{\mathbb{Z}} \mathbb{Q} \longrightarrow
\mathbb{Q}[H]^{\oplus 9} \longrightarrow I[H] \otimes_{\mathbb{Z}} \mathbb{Q}
\longrightarrow 0 .
\]
The augmentation sequence
\begin{equation}\label{eq:aug}
0 \longrightarrow I[H] \longrightarrow \mathbb{Z}[H]
\xrightarrow{\ \varepsilon\ } \mathbb{Z} \longrightarrow 0
\end{equation}
remains exact after tensoring with \(\mathbb{Q}\) (flat over
\(\mathbb{Z}\), \cite[Ch.II, \S 2, no.4, Th.1]{Bourbaki2006}), giving
\[
0 \longrightarrow I[H] \otimes_{\mathbb{Z}} \mathbb{Q} \longrightarrow
\mathbb{Q}[H] \xrightarrow{\ \varepsilon\ } \mathbb{Q} \longrightarrow 0 .
\]
This sequence splits since \(|H| = 9\) is invertible in
\(\mathbb{Q}\) (Maschke's Theorem \cite[(3.14)]{CurtisReiner1981};
the splitting is given by the idempotent
\(e_0 = \frac{1}{|H|}\sum_{h \in H} h\)).  Hence in the Grothendieck group
\(K_0(\mathbb{Q}[H])\) \cite[(16.3)]{CurtisReiner1981},
\[
[I[H] \otimes_{\mathbb{Z}} \mathbb{Q}] = [\mathbb{Q}[H]] - [\mathbb{Q}].
\]
From the exact sequence of \(\mathbb{Q}[H]\)-modules above,
\begin{align*}
[M|_H \otimes \mathbb{Q}]
&= 9[\mathbb{Q}[H]] - [I[H] \otimes \mathbb{Q}] \\
&= 9[\mathbb{Q}[H]] - ([\mathbb{Q}[H]] - [\mathbb{Q}]) \\
&= 8[\mathbb{Q}[H]] + [\mathbb{Q}].
\end{align*}
Since \(\mathbb{Q}[H]\) is semisimple by Maschke
\cite[(3.14)]{CurtisReiner1981}, equality in \(K_0(\mathbb{Q}[H])\) implies
\[
M|_H \otimes_{\mathbb{Z}} \mathbb{Q} \cong 8\cdot \mathbb{Q}[H] \oplus
\mathbb{Q}
\]
of \(\mathbb{Q}[H]\)-modules.
\end{proof}

% \begin{proposition}\label{prop:QH-decomp}
% \[
% \mathbb{Q}[H] \cong \mathbb{Q} \times \mathbb{Q}(\zeta_3)^4 ,
% \]
% as \(\mathbb{Q}\)-algebras, where \(\zeta_3\) denotes a primitive third
% root of unity.
% \end{proposition}
%
% \begin{proof}
% By the classification of rational group algebras of finite abelian
% groups \cite[\S10B]{CurtisReiner1981}: the characters of
% \(H = C_3 \times C_3\) form four Galois orbits of size~\(2\)
% (the \(8\) non-trivial characters, paired by complex conjugation
% under \(\Gal(\mathbb{Q}(\zeta_3)/\mathbb{Q}) \cong C_2\))
% and one of size~\(1\) (the trivial character).
% \end{proof}

From Proposition~\ref{prop:MH-char} we deduce: if \(M|_H\) \emph{were} a
permutation \(\mathbb{Z}[H]\)-lattice, its rational character would force \(M|_H
\cong 8\cdot\mathbb{Z}[H] \oplus \mathbb{Z}\).

\begin{lemma}\label{lem:perm-class}
Let \(P\) be a permutation \(\mathbb{Z}[H]\)-lattice satisfying
\(P \otimes_{\mathbb{Z}} \mathbb{Q} \cong 8\cdot\mathbb{Q}[H] \oplus
\mathbb{Q}\). Then \(P \cong 8\cdot\mathbb{Z}[H] \oplus \mathbb{Z}\).
\end{lemma}

\begin{proof}
Since \(P\) is a permutation \(\mathbb{Z}[H]\)-lattice
(Definition~\ref{def:lattice}), we can write
\[
P \cong \bigoplus_{K \le H} m_K \cdot \mathbb{Z}[H/K], \qquad
m_K \in \mathbb{Z}_{\ge 0},
\]
where \(K\) runs over the subgroups of \(H = C_3 \times C_3\). The subgroups
are: \(\{1\}\) (trivial, index \(9\)); four subgroups of order \(3\) (each index
\(3\)), namely \(K_1 = \langle a \rangle\), \(K_2 = \langle b \rangle\), \(K_3 =
\langle ab \rangle\), \(K_4 = \langle ab^2 \rangle\) where \(H = \langle a
\rangle \times \langle b \rangle \cong C_3 \times C_3\); and \(H\) itself (index
\(1\), giving the trivial module \(\mathbb{Z}\)).

\medskip\noindent
\emph{Character-theoretic determination of multiplicities.}
By Lemma~\ref{lem:perm-char}, the character of \(\mathbb{Q}[H/K]\)
is given by
\[
\chi_{H/K}(g) = \begin{cases}
[H:K], & g \in K,\\
0, & g \notin K,
\end{cases}
\]
and evaluating the character of \(P\) at \(g \in H\) gives
\[
\chi_P(g) = \sum_{K \ni g} m_K\, [H:K].
\]
% [2026-10-01 第二十改] 原句「and the multiplicities \(m_K\) are uniquely
% determined by the character of~\(P\).」替换为上求值式（该唯一性断言为假；
% 原句及依据存档见 lem:perm-char 之后注）。

By Proposition~\ref{prop:MH-char},
\(\chi_P = 8\chi_{H/\{1\}} + \chi_{H/H}\),
so \(\chi_P(1) = 8 \cdot 9 + 1 = 73\) and
\(\chi_P(g) = 0 + 1 = 1\) for all \(g \neq 1\).

\medskip\noindent
\emph{System of equations.}
Let \(a = m_{\{1\}}\), \(b_j = m_{K_j}\) (\(j=1,\dots,4\)), and \(c = m_H\).
Evaluating \(\chi_P\) at the identity and at
\(g \in K_j \setminus \{1\}\) gives
\begin{align}
\chi_P(1) &= 9a + 3(b_1+b_2+b_3+b_4) + c = 73,
\label{eq:sys-id}\\
\chi_P(g) &= 3b_j + c = 1 \qquad (g \in K_j,\ g \neq 1,\ j=1,\dots,4).
\label{eq:sys-order3}
\end{align}

\medskip\noindent
\emph{Solution.}
Equation~\eqref{eq:sys-order3} gives \(b_j = (1-c)/3\). Since the multiplicities
\(b_j\) are non-negative integers by definition of a permutation lattice
decomposition, we must have \(c = 1\), forcing \(b_j = 0\) for all \(j\).
Substituting into~\eqref{eq:sys-id}: \(9a + 0 + 1 = 73\), hence \(a = 8\).
Therefore \(P \cong 8\cdot\mathbb{Z}[H] \oplus \mathbb{Z}\).
\end{proof}

% \begin{remark}\label{rem:integrality} The argument above is entirely
% elementary: the multiplicities \(m_K\) are integers by the definition of a
% permutation \(\mathbb{Z}[H]\)-lattice, and the character evaluation at
% specific group elements yields an invertible integer-linear system that
% determines them uniquely. No passage through the Burnside ring or delicate
% integrality theorems is required. \end{remark}

%%%%%%%%%%%%%%%%%%%%%%%%%%%%%%%%%%%%%%%%%%%%%%%%%%%%%%%%%%%%%%%%%%%%%%%%%%%
\section{Cohomological Obstruction to Stable Permutation}
\label{sec:cohom}
%%%%%%%%%%%%%%%%%%%%%%%%%%%%%%%%%%%%%%%%%%%%%%%%%%%%%%%%%%%%%%%%%%%%%%%%%%%

In this section we compute the second integral cohomology group
\(\coh^2(H, M|_H)\) and prove that its exponent obstructs stable permutation.

\medskip\noindent
\textbf{Notation.}
Throughout this section we write \(\coh^q(G, M)\) for the \(q\)-th
integral group cohomology of a \(G\)-module \(M\).  The upright~\(\coh\)
is reserved for the cohomology functor, while the italic~\(H\) denotes
the group \(H = C_3 \times C_3\).

%%%%%%%%%%%%%%%%%%%%%%%%%%%%%%%%%%%%%%%%%%%%%%%%%%%%%%%%%%%%%%%%%%%%%%%%%%%
\subsection{Cohomology Computations}

\begin{lemma}\label{lem:H1-IH}
The group \(\coh^1(H, I[H])\) is cyclic of order \(9\).
\end{lemma}

\begin{proof}
Apply group cohomology \(\coh^*(H, -)\) to the augmentation exact sequence
\eqref{eq:aug}.
The associated long exact cohomology sequence reads
\begin{multline*}
0 \longrightarrow \coh^0(H, I[H]) \longrightarrow \coh^0(H, \mathbb{Z}[H])
\longrightarrow \coh^0(H, \mathbb{Z}) \\
\longrightarrow \coh^1(H, I[H])
\longrightarrow \coh^1(H, \mathbb{Z}[H]) \longrightarrow \cdots .
\end{multline*}
We analyse the low-degree terms:
\begin{itemize}
\item \(\coh^0(H, \mathbb{Z}[H]) = \mathbb{Z}[H]^H = \mathbb{Z}\!\cdot\! N_H\),
  where \(N_H = \sum_{h \in H} h\) is the norm element.
\item \(\coh^0(H, \mathbb{Z}) = \mathbb{Z}^H = \mathbb{Z}\).
\item The map \(\coh^0(H, \mathbb{Z}[H]) \to \coh^0(H, \mathbb{Z})\) sends
  \(N_H \mapsto \varepsilon(N_H) = \sum_{h \in H} 1 = |H| = 9\).
\item \(\coh^1(H, \mathbb{Z}[H]) = 0\): \(\mathbb{Z}[H] \cong \Ind_1^H
\mathbb{Z}\)
  as \(\mathbb{Z}[H]\)-modules, and Shapiro's Lemma
  \cite[Theorem~6.3.2]{Weibel1994} gives
  \(\coh^1(H, \Ind_1^H \mathbb{Z}) \cong \coh^1(\{1\}, \mathbb{Z}) = 0\).
\end{itemize}
From the exactness at \(\coh^0(H, \mathbb{Z})\) we obtain
\(\coker(\mathbb{Z}\!\cdot\!N_H \xrightarrow{\varepsilon} \mathbb{Z})
= \mathbb{Z}/9\mathbb{Z}\), and exactness at \(\coh^1(H, I[H])\) with the
vanishing of \(\coh^1(H, \mathbb{Z}[H])\) yields the short exact sequence
\[
0 \longrightarrow \mathbb{Z}/9\mathbb{Z} \longrightarrow \coh^1(H, I[H])
\longrightarrow 0,
\]
hence \(\coh^1(H, I[H]) \cong \mathbb{Z}/9\mathbb{Z} \cong C_9\).
\end{proof}

\begin{lemma}\label{lem:H2-MH}
For the Procesi restriction \(M|_H\),
\(\coh^2(H, M|_H) \cong C_9\).
\end{lemma}

\begin{proof}
Apply \(\coh^*(H,-)\) to the exact sequence of
Proposition~\ref{prop:procesi-H}:
\[
0 \longrightarrow M|_H \longrightarrow \mathbb{Z}[H]^{\oplus 9}
\longrightarrow I[H] \longrightarrow 0 .
\]
The associated long exact cohomology sequence gives
\[
\begin{aligned}
\cdots &\longrightarrow \coh^1(H, \mathbb{Z}[H]^{\oplus 9})
\longrightarrow \coh^1(H, I[H]) \\
&\longrightarrow \coh^2(H, M|_H)
\longrightarrow \coh^2(H, \mathbb{Z}[H]^{\oplus 9}) \longrightarrow \cdots .
\end{aligned}
\]

By Shapiro's Lemma \cite[Theorem~6.3.2]{Weibel1994},
\(\coh^i(H, \mathbb{Z}[H]) \cong \coh^i(\{1\}, \mathbb{Z}) = 0\) for all \(i >
0\). Hence \(\coh^1(H, \mathbb{Z}[H]^{\oplus 9}) = \coh^2(H,
\mathbb{Z}[H]^{\oplus 9}) = 0\), and the long exact sequence collapses to the
isomorphism
\[
\coh^2(H, M|_H) \;\cong\; \coh^1(H, I[H]).
\]
By Lemma~\ref{lem:H1-IH}, \(\coh^1(H, I[H]) \cong C_9\).
\end{proof}

\begin{lemma}\label{lem:H2-Z}
With the trivial \(H\)-action on \(\mathbb{Z}\),
\(\coh^2(H, \mathbb{Z}) \cong C_3 \times C_3\).
\end{lemma}

\begin{proof}
Apply \(\coh^*(H, -)\) to the short exact sequence of trivial \(H\)-modules
\[
0 \longrightarrow \mathbb{Z} \longrightarrow \mathbb{Q}
\longrightarrow \mathbb{Q}/\mathbb{Z} \longrightarrow 0 .
\]
The associated long exact sequence reads
\begin{multline*}
0 \longrightarrow \mathbb{Z} \longrightarrow \mathbb{Q}
\longrightarrow \mathbb{Q}/\mathbb{Z} \\
\longrightarrow \coh^1(H, \mathbb{Z})
\longrightarrow \coh^1(H, \mathbb{Q})
\longrightarrow \coh^1(H, \mathbb{Q}/\mathbb{Z}) \\
\longrightarrow \coh^2(H, \mathbb{Z})
\longrightarrow \coh^2(H, \mathbb{Q}) \longrightarrow \cdots .
\end{multline*}
Since \(\mathbb{Q}\) is uniquely divisible, multiplication by \(3\) is an
automorphism of \(\mathbb{Q}\); hence \(\coh^i(H, \mathbb{Q}) = 0\) for all
\(i > 0\) \cite[Proposition~1.6.2(c)]{NSW2020}.  Substituting
\(\coh^1(H, \mathbb{Q}) = \coh^2(H, \mathbb{Q}) = 0\) yields
\(\coh^2(H, \mathbb{Z}) \cong \coh^1(H, \mathbb{Q}/\mathbb{Z})\).
Since \(\mathbb{Q}/\mathbb{Z}\) carries the trivial \(H\)-action,
\(\coh^1(H, \mathbb{Q}/\mathbb{Z}) = \Hom(H, \mathbb{Q}/\mathbb{Z})\)
\cite[Corollary~6.4.6]{Weibel1994}.
Now \(\Hom(H, \mathbb{Q}/\mathbb{Z}) \cong C_3 \times C_3\) because
\(\Hom(C_3, \mathbb{Q}/\mathbb{Z}) \cong C_3\)
(\(\mathbb{Q}/\mathbb{Z}\) has a unique subgroup of order~\(3\))
and \(\Hom({-}, \mathbb{Q}/\mathbb{Z})\) preserves finite direct sums
\cite[Ch.II, \S1, no.6, Prop.6, Cor.1]{Bourbaki2007a}.
Hence \(\coh^2(H, \mathbb{Z}) \cong C_3 \times C_3\).
\end{proof}

%%%%%%%%%%%%%%%%%%%%%%%%%%%%%%%%%%%%%%%%%%%%%%%%%%%%%%%%%%%%%%%%%%%%%%%%%%%
\subsection{Exponent of \(\coh^2\) for Permutation Lattices}

\begin{proposition}\label{prop:perm-exponent}
Let \(P\) be a permutation \(\mathbb{Z}[H]\)-lattice.  Then every element of
\(\coh^2(H, P)\) has order dividing \(3\); equivalently,
\(\exp(\coh^2(H, P)) \mid 3\).
\end{proposition}

\begin{proof}
Write \(P \cong \bigoplus_{i} \mathbb{Z}[H/K_i]\) for subgroups \(K_i \le H\).
Since \(\coh^2(H, {-})\) is additive and therefore preserves finite
direct sums \cite[Proposition~3.3.4]{Weibel1994},
\(\coh^2(H, P) \cong \bigoplus_i \coh^2(H, \mathbb{Z}[H/K_i])\).
By Shapiro's Lemma \cite[Theorem~6.3.2]{Weibel1994},
\(\mathbb{Z}[H/K_i] \cong \Ind_{K_i}^H(\mathbb{Z})\),
hence \(\coh^2(H, \mathbb{Z}[H/K_i]) \cong \coh^2(K_i, \mathbb{Z})\),
where \(\mathbb{Z}\) carries the trivial \(K_i\)-action.
For each subgroup \(K \le H\) we compute \(\coh^2(K, \mathbb{Z})\):
\begin{itemize}
  \item \(K = \{1\}\):
    \(\coh^2(\{1\}, \mathbb{Z}) = 0\).
  \item \(K \cong C_3\) (four subgroups of order \(3\)):
\(\coh^2(C_3, \mathbb{Z}) \cong \Hom(C_3, \mathbb{Q}/\mathbb{Z}) \cong C_3\).
  \item \(K = H = C_3 \times C_3\):
    \(\coh^2(H, \mathbb{Z}) \cong C_3 \times C_3\) (Lemma~\ref{lem:H2-Z}).
\end{itemize}
In every case \(\exp(\coh^2(K, \mathbb{Z})) \mid 3\).  Hence
\(\exp(\coh^2(H, P)) \mid 3\).
\end{proof}

\begin{proposition}\label{prop:stable-exp}
If a \(\mathbb{Z}[H]\)-lattice \(L\) is stably permutation, then
\[ \exp(\coh^2(H, L)) \mid 3. \]
\end{proposition}
\begin{proof}
Write \(L \oplus P \cong Q\) with \(P, Q\) permutation.  By additivity
\cite[Proposition~3.3.4]{Weibel1994},
\(\coh^2(H, L) \oplus \coh^2(H, P) \cong \coh^2(H, Q)\).  Both
\(\coh^2(H, P)\) and \(\coh^2(H, Q)\) have exponent dividing \(3\) by
Proposition~\ref{prop:perm-exponent}.  For any \(x \in \coh^2(H, L)\),
the element \((x,0) \in \coh^2(H, Q)\) satisfies \(3(x,0) = 0\); hence
\(3x = 0\).  Therefore \(\exp(\coh^2(H, L)) \mid 3\).
\end{proof}

%%%%%%%%%%%%%%%%%%%%%%%%%%%%%%%%%%%%%%%%%%%%%%%%%%%%%%%%%%%%%%%%%%%%%%%%%%%
\subsection{The Obstruction}

\begin{theorem}\label{thm:not-stably-perm}
The restricted Procesi lattice \(M|_H\) is not a stably permutation
\(\mathbb{Z}[H]\)-lattice.
\end{theorem}

\begin{proof}
By Lemma~\ref{lem:H2-MH}, \(\coh^2(H, M|_H) \cong C_9\), a group of exponent
\(9\).  If \(M|_H\) were stably permutation, Proposition~\ref{prop:stable-exp}
would force \(\exp(\coh^2(H, M|_H)) \mid 3\).  Contradiction.
\end{proof}

\begin{corollary}\label{thm:not-stably-rational}
The field \(Z_H(F,9)\) is not stably rational over \(F\).
\end{corollary}

\begin{proof}
Recall that \(Z_H(F,9) \cong F(M|_H)^H\)
(Remark~\ref{rem:procesi-geometric}).
By Theorem~\ref{thm:not-stably-perm}, \(M|_H\) is not stably permutation.
By Theorem~\ref{thm:endo-miyata}, used as a hypothesis,
stable rationality of \(F(M|_H)^H\) over \(F\) forces \(M|_H\) to be
stably permutation.  Hence \(Z_H(F,9)\) is not stably rational.
\end{proof}

\begin{remark}
% The proof uses only basic homological algebra (Shapiro's lemma,
% cohomology long exact sequences, additivity of \(\coh^2\)) and the
% elementary fact that \(C_9\) has exponent \(9\) while permutation
% lattices have \(\coh^2\) of exponent dividing \(3\).  No \(3\)-adic
% completions, maximal orders, Swan subgroups, or flasque class group
% computations are required.
The purely lattice-theoretic part (Theorem~\ref{thm:not-stably-perm})
holds integrally, hence for any field whose characteristic does not
divide \(|H| = 9\).  The geometric conclusion
(Theorem~\ref{thm:not-stably-rational}) additionally requires
characteristic~\(0\) via Theorem~\ref{thm:endo-miyata}.

\begin{proposition}\label{prop:faithful-MH}
Let \(M_9(F)\) denote the algebra of \(9 \times 9\) matrices over \(F\).
Let \(H = C_3 \times C_3\) act on \(M_9(F)\) via the regular representation
\(H \hookrightarrow S_9\) by permutation of matrix indices:
\(h \cdot e_{ij} = e_{h(i), h(j)}\) for the standard matrix units.
The action of \(H\) on \(M|_H\) is the restriction of the diagonal
left-regular action on \(\mathbb{Z}[H]^{\oplus 9}\) of
Proposition~\ref{prop:tensor-reg-ZH}:
\[
h \cdot (z_b)_{b \in H} = (h\, z_{h^{-1} b'})_{b' \in H},
\qquad
h \in H,\; (z_b)_{b \in H} \in M|_H \subset \mathbb{Z}[H]^{\oplus 9},
\]
that is, left multiplication composed with the coordinate permutation
\(b \mapsto h b\).
This action is faithful.
\end{proposition}

\begin{proof}
The action on \(M_9(F)\) is faithful because \(h \neq 1\) implies
\(h(i) \neq i\) for some index~\(i\), whence \(h \cdot e_{ij} \neq e_{ij}\).
The Procesi isomorphism \(Z_H(F,9) \cong F(M|_H)^H\)
(Remark~\ref{rem:procesi-geometric}) intertwines the \(H\)-actions.
Any \(h \in H\) acting trivially on \(M|_H\) would, by functoriality of
the field of fractions, act trivially on
\(F(M|_H) = \operatorname{Frac}(\operatorname{Sym}_F(M|_H
\otimes_{\mathbb{Z}} F))\), where \(\operatorname{Sym}_F\) denotes the
symmetric algebra over \(F\), and therefore on
\(F(M|_H)^H \cong Z_H(F,9)\), hence on \(M_9(F)\) ---
contradicting faithfulness on \(M_9(F)\).
\end{proof}

\begin{proof}[Second proof]
Let \(h \in H\), \(h \neq 1\), and let \(z \in
\mathbb{Z}[H]^{\oplus 9}\) be the element with coordinates
\(z_1 = h - 1\), \(z_h = 1 - h\), and \(z_b = 0\) for
\(b \notin \{1, h\}\).  Under the identifications of
Proposition~\ref{prop:tensor-reg-ZH}, the element with coordinates
\((z_b)_{b \in H}\) corresponds to \(\sum_{b \in H} z_b \otimes
\delta_b\), where \(\delta_b\) denotes the group element \(b\) of the
second factor, identified with \(e_{b \cdot \omega_0}\).  Since
\(\varphi(e_\alpha \otimes e_\beta) = e_\alpha - e_\beta\), the map
\(\pi|_H\) sends this element to
\(\sum_{b \in H}(z_b - \varepsilon(z_b)\, b)\).  Since
\(\varepsilon(h - 1) = \varepsilon(1 - h) = 0\), this gives
\(\pi|_H(z) = (h - 1) + (1 - h) = 0\), and \(z \in M|_H\).
Now act by \(h\).  By Proposition~\ref{prop:tensor-reg-ZH} the
\(h^{-1}\)-th coordinate of \(z\) moves into the first coordinate and
is multiplied by \(h\), so \((h \cdot z)_1 = h \cdot z_{h^{-1}}\).
Here \(h^{-1} \neq 1\) because \(h \neq 1\), and \(h^{-1} \neq h\):
if \(h^{-1} = h\), then \(h^2 = 1\), and \(h = h^3 (h^2)^{-1} = 1\)
since \(h^3 = 1\).  Hence \(z_{h^{-1}} = 0\), and \((h \cdot z)_1 = 0\),
whereas \(z_1 = h - 1 \neq 0\).  Therefore \(h \cdot z \neq z\), that
is, \(h\) acts nontrivially on \(M|_H\).
\end{proof}

\begin{theorem}\label{thm:saltman-cor3.13}
\cite[Corollary~3.13]{Saltman1984}.
Let \(A\) be a finite abelian group, and let \(2^r\) be the highest power
of \(2\) dividing the exponent of \(A\).  Assume \(F\) is a field such that
either \(F\) has characteristic \(2\) or \(F(\rho)/F\) is a cyclic extension,
where \(\rho\) is a primitive \(2^r\)-th root of unity.  Let \(V\) be a
finitely generated \(F[A]\)-module on which \(A\) acts faithfully.  Then
\(F(V)^A\) is retract rational over \(F\).
\end{theorem}

\begin{corollary}\label{thm:retract}
The field \(Z_H(F,9)\) is retract rational over \(F\).
% [2026-10-01 去重] 原句尾括号 (Remark~\ref{rem:procesi-geometric}) 已删——
% 证明首行已引同一 remark，陈述不再重复。
\end{corollary}

\begin{proof}
\(Z_H(F,9) \cong F(M|_H)^H\) (Remark~\ref{rem:procesi-geometric}).
Let \(V = M|_H \otimes_{\mathbb{Z}} F\); then \(F(V)^H \cong F(M|_H)^H\).
Apply Theorem~\ref{thm:saltman-cor3.13} with \(A = H = C_3 \times C_3\).
The exponent of \(H\) is \(3\) (odd), so \(2^r = 1\) and \(\rho = 1\); the
condition on \(F(1)/F\) is vacuous.  The group \(H\) acts faithfully on \(V\)
because it acts faithfully on \(M|_H\) (Proposition~\ref{prop:faithful-MH}) and
\(\operatorname{char} F = 0\).  Hence \(Z_H(F,9)\) is retract rational
over \(F\).
\end{proof}
\end{remark}

%%%%%%%%%%%%%%%%%%%%%%%%%%%%%%%%%%%%%%%%%%%%%%%%%%%%%%%%%%%%%%%%%%%%%%%%%%%
\begin{proposition}\label{prop:M_H-not-permutation}
The restricted Procesi lattice \(M|_H\) is not a permutation
\(\mathbb{Z}[H]\)-lattice.
\end{proposition}
\begin{proof}
Apply \(\Hom_{\mathbb{Z}[H]}(-, L)\) to \eqref{eq:aug}.  The associated
long exact sequence reads
\begin{multline*}
0 \longrightarrow \Hom_{\mathbb{Z}[H]}(\mathbb{Z}, L)
\longrightarrow \Hom_{\mathbb{Z}[H]}(\mathbb{Z}[H], L)
\longrightarrow \Hom_{\mathbb{Z}[H]}(I[H], L) \\
\longrightarrow \Ext^1_{\mathbb{Z}[H]}(\mathbb{Z}, L)
\longrightarrow \Ext^1_{\mathbb{Z}[H]}(\mathbb{Z}[H], L)
\longrightarrow \Ext^1_{\mathbb{Z}[H]}(I[H], L) \\
\longrightarrow \Ext^2_{\mathbb{Z}[H]}(\mathbb{Z}, L)
\longrightarrow \Ext^2_{\mathbb{Z}[H]}(\mathbb{Z}[H], L)
\longrightarrow \cdots .
\end{multline*}
Since \(\mathbb{Z}[H]\) is projective,
\(\Ext^i_{\mathbb{Z}[H]}(\mathbb{Z}[H], {-}) = 0\) for all \(i \ge 1\).
Substituting this collapses the sequence to an isomorphism
\(\Ext^1_{\mathbb{Z}[H]}(I[H], L) \cong \Ext^2_{\mathbb{Z}[H]}(\mathbb{Z}, L)
= \coh^2(H, L)\).
Taking \(L = M|_H\) and \(L = 8\cdot\mathbb{Z}[H] \oplus \mathbb{Z}\)
yields
\[
\Ext^1_{\mathbb{Z}[H]}(I[H], M|_H) \cong C_9,\qquad
\Ext^1_{\mathbb{Z}[H]}(I[H], 8\cdot\mathbb{Z}[H] \oplus \mathbb{Z})
\cong C_3 \times C_3 .
\]
By Proposition~\ref{prop:MH-char},
\(M|_H \otimes_{\mathbb{Z}} \mathbb{Q} \cong
8\cdot\mathbb{Q}[H] \oplus \mathbb{Q}\).  If \(M|_H\) were a
permutation lattice, Lemma~\ref{lem:perm-class} would then give
\(M|_H \cong 8\cdot\mathbb{Z}[H] \oplus \mathbb{Z}\), and the two groups
above would coincide, contradicting \(C_9 \not\cong C_3 \times C_3\).
\end{proof}

\begin{proof}[Second proof]
By Lemma~\ref{lem:H2-MH}, \(\coh^2(H, M|_H) \cong C_9\), a group of
exponent \(9\).  If \(M|_H\) were a permutation
\(\mathbb{Z}[H]\)-lattice, Proposition~\ref{prop:perm-exponent} would
force \(\exp(\coh^2(H, M|_H)) \mid 3\).  Contradiction.
\end{proof}

\begin{remark}
Proposition~\ref{prop:M_H-not-permutation} is strictly weaker than
Theorem~\ref{thm:not-stably-perm} and cannot replace it in the proof of
Corollary~\ref{thm:not-stably-rational}.  Indeed a permutation lattice
is stably permutation (take \(P = 0\)), so non-permutation does not
imply non-stable permutation; the main proof therefore needs the
stronger statement of Theorem~\ref{thm:not-stably-perm} together with
the implication ``stably rational \(\Rightarrow\) stably permutation''
(Theorem~\ref{thm:endo-miyata}).
\end{remark}

\section{Generalization to \(C_p \times C_p\)}
\label{sec:generalization}
%%%%%%%%%%%%%%%%%%%%%%%%%%%%%%%%%%%%%%%%%%%%%%%%%%%%%%%%%%%%%%%%%%%%%%%%%%%

We record the generalization of Theorem~\ref{thm:not-stably-rational} to all
odd primes.  The comparison with the unramified Brauer group, and the scope of
the regular embedding used above, are deferred to
Appendix~\ref{app:brauer} and Appendix~\ref{app:scope}.

\begin{theorem}\label{thm:Cp-general}
Let \(p\) be an odd prime and let \(H = C_p \times C_p \subset S_{p^2}\)
act regularly (i.e., freely and transitively on
\(\{1, 2, \dots, p^2\}\)).  Let \(M|_H\) be the restriction of the Procesi
\(\mathbb{Z}[S_{p^2}]\)-lattice to \(H\).  By the regular identification
\(U|_H \cong \mathbb{Z}[H]\) and the projection formula
\(\mathbb{Z}[H] \otimes_{\mathbb{Z}} \mathbb{Z}[H] \cong
\mathbb{Z}[H]^{\oplus p^2}\),
% [2026-10-01 去重] 原括号引用 (Proposition~\ref{prop:tensor-reg-ZH}) 已删——
% 证明首段已引同一命题，陈述不再重复。
\(M|_H\) fits into the exact sequence of \(\mathbb{Z}[H]\)-modules
\[
0 \longrightarrow M|_H \longrightarrow \mathbb{Z}[H]^{\oplus p^2}
\longrightarrow I[H] \longrightarrow 0 ,
\]
where \(I[H] = \ker(\varepsilon \colon \mathbb{Z}[H] \twoheadrightarrow
\mathbb{Z})\) is the augmentation ideal and the surjection sends each
standard basis element of \(\mathbb{Z}[H]^{\oplus p^2}\) to a difference
\(h - h'\) in \(\mathbb{Z}[H]\).
Let \(F\) be a field of characteristic \(0\), and let
\(Z_H(F, p^2)\) denote the center of the generic division algebra of
degree~\(p^2\) over \(F\) (cf.\ Remark~\ref{rem:procesi-geometric}),
with Procesi isomorphism \(Z_H(F, p^2) \cong F(M|_H)^H\).  Then:
\begin{enumerate}
  \item \(\coh^2(H, M|_H) \cong C_{p^2}\).
  \item For any permutation \(\mathbb{Z}[H]\)-lattice \(P\),
    \(\exp(\coh^2(H, P)) \mid p\).
  \item Consequently, \(M|_H\) is not stably permutation, and
    \(Z_H(F, p^2)\) is not stably rational (hence not rational) over \(F\).
  \item \(Z_H(F, p^2)\) is retract rational over \(F\).
\end{enumerate}
\end{theorem}

\begin{proof}
The Procesi restriction to \(H\) follows the same construction as
in \S\ref{sec:procesi}: the regular identification \(U|_H \cong \mathbb{Z}[H]\)
(Proposition~\ref{prop:U-restrict-canonical})
and the isomorphism
\(\mathbb{Z}[H] \otimes \mathbb{Z}[H] \cong \mathbb{Z}[H]^{\oplus p^2}\)
(Proposition~\ref{prop:tensor-reg-ZH}, whose proof works verbatim
with~\(9\) replaced by~\(p^2\)) hold integrally for
all~\(p\).  The map \(\pi|_H : \mathbb{Z}[H]^{\oplus p^2} \to I[H]\) is
surjective because \(I[H]\) is generated by \(h-1\) (\(h \neq 1\)), and
\(M|_H = \ker(\pi|_H)\) is \(\mathbb{Z}\)-free (submodule of a free
\(\mathbb{Z}\)-module) of \(\mathbb{Z}\)-rank \(p^4 - p^2 + 1\).

\medskip\noindent
(1)~The same cohomological computation as in Lemmas~\ref{lem:H1-IH}
and~\ref{lem:H2-MH}, with \(9\) replaced by \(p^2\) and \(H = C_p \times C_p\),
gives \(\coh^2(H, M|_H) \cong C_{p^2}\).

\medskip\noindent
(2)~The same argument as in Proposition~\ref{prop:perm-exponent},
with the subgroups of \(C_p \times C_p\) in place of those of
\(C_3 \times C_3\), gives \(\exp(\coh^2(H, P)) \mid p\).

\medskip\noindent
(3)~The same exponent argument as in Theorem~\ref{thm:not-stably-perm}
and Corollary~\ref{thm:not-stably-rational}, with \(p\) in place of~\(3\),
shows that \(M|_H\) is not stably permutation and
\(Z_H(F, p^2)\) is not stably rational.

\medskip\noindent
(4)~The same argument as in Theorem~\ref{thm:retract} applies:
\(H\) is abelian of odd exponent, so Saltman's
\cite[Corollary~3.13]{Saltman1984} gives retract rationality.
\end{proof}

\begin{remark}
The proof of Theorem~\ref{thm:Cp-general} holds for all odd primes~\(p\)
without any dependence on~\(p=3\) specifically:
the regular embedding \(H \hookrightarrow S_{p^2}\)
(independent of the choice up to conjugation;
Proposition~\ref{prop:embedding-independence}),
Proposition~\ref{prop:tensor-reg-ZH},
the cohomological vanishing (Shapiro), and the subgroup
analysis of \(C_p \times C_p\) are all purely formal and independent
of~\(p\).  The converse
Endo--Miyata implication (rational \(\Rightarrow\) stably permutation for
non-cyclic \(p\)-groups) remains open but is not required for any of the
above conclusions.
\end{remark}

Our results on \cite[Problem~5.2]{ABGV2011} are summarized in
Table~\ref{tab:main}.

% The central new result is \(\coh^2(H, M|_H) \cong C_9\)
% (Lemma~\ref{lem:H2-MH}) and the exponent argument showing that this obstructs
% stable permutation (Theorem~\ref{thm:not-stably-perm}). The chain of
% implications \[ \begin{aligned} &\coh^2 \cong C_9 \\ &\;\Longrightarrow\;
% \text{not stably permutation} \qquad(\exp \mid 3 \text{ for permutation
% lattices}) \\ &\;\Longrightarrow\; \text{not stably rational}
% \qquad(\text{Theorem~\ref{thm:endo-miyata}}) \\ &\;\Longrightarrow\; \text{not
% rational} \qquad(\text{rational} \Rightarrow \text{stably rational})
% \end{aligned} \] resolves both the stable rationality and rationality
% questions simultaneously (in characteristic~\(0\), without reliance on the
% open converse Endo--Miyata implication). The \(\Ext^1\) computation
% (Proposition~\ref{prop:M_H-not-permutation}) provides independent confirmation
% that \(M|_H\) is not a permutation lattice. The Bogomolov multiplier
% computation \(B_0(C_3\times C_3) \cong C_3\) (Proposition~\ref{prop:B0}) shows
% that the classical unramified Brauer group of the Noether problem \(F(H)^H\)
% is non-zero, while the integral \(\coh^2\) obstruction operates at the level
% of \(\mathbb{Z}[H]\)-lattices and is of a different type. \medskip\noindent
% \textbf{Scope, limitations, and further
% directions.}\label{sec:scope-limitations} \textbf{Generalization to \(C_p
% \times C_p\).} The entire proof generalizes to all odd primes~\(p\) (the
% cohomological part is purely integral; the geometric conclusion requires
% \(\Char F = 0\)) (see \S\ref{sec:generalization-roadmap} above). For \(p=5\),
% a direct computational verification (Appendix~\ref{sec:gap-p5}) confirms the
% pattern: \(\coh^1(H, \mathbb{Z}[H])=0\) (free rank~0,
% \(\operatorname{rank}(d^0)=24\), \(\operatorname{rank}(d^1)=601\)),
% \(\coh^1(H, I[H]) \cong C_{25}\) (free rank~0,
% \(\operatorname{rank}(d^0)=24\), \(\operatorname{rank}(d^1)=576\)),
% \(\coh^2(H, \mathbb{Z}) \cong C_5 \times C_5\), and consequently \(\coh^2(H,
% M|_H) \cong C_{25}\) with exponent~\(25 \nmid 5\). All steps are purely formal
% and independent of~\(p\); no class-group or \(p\)-adic integrality conditions
% intervene. The converse Endo--Miyata implication remains open for all
% non-cyclic \(p\)-groups, but is not required for the main results (since
% \(\text{not stably perm} \;\Longrightarrow\; \text{not stably rational}
% \;\Longrightarrow\; \text{not rational}\)).
% [2026-10-01 用户指令] 附录 A 节标题中「Computational Verification」缩为「Computation」。
\section{Appendix A: Explicit Cocycle and Computation}
\label{app:cocycle}
%%%%%%%%%%%%%%%%%%%%%%%%%%%%%%%%%%%%%%%%%%%%%%%%%%%%%%%%%%%%%%%%%%%%%%%%%%%

% [2026-10-01 用户指令] 附录 A 子节标题降级为段标题（\subsection → \paragraph；
% 该标题无 live \ref）
\paragraph{Explicit 2-Cocycle Representative}

We construct an explicit \(2\)-cocycle \(z : H \times H \to M|_H\) whose
cohomology class generates \(\coh^2(H, M|_H) \cong C_9\).

The connecting homomorphism \(\partial : \coh^1(H, I[H]) \xrightarrow{\cong}
\coh^2(H, M|_H)\) from the Procesi sequence sends a \(1\)-cocycle \(\psi\) to
\(z(g,h) = g \cdot \widetilde{\psi}(h) - \widetilde{\psi}(gh) +
\widetilde{\psi}(g)\), where \(\widetilde{\psi} : H \to \mathbb{Z}[H]^{\oplus
9}\) is any set-theoretic lift of \(\psi\).

\medskip\noindent
\textbf{Step 1: generator of \(\coh^1(H, I[H])\).}
\(\psi(h) = h-1 \in I[H]\) (Lemma~\ref{lem:H1-IH}).

\medskip\noindent
\textbf{Step 2: splitting of \(\pi\).}
Recall \(\pi(x_1,\dots,x_9) = \sum_{h \in H}(x_h - \varepsilon(x_h)h)\).
Define \(s : I[H] \to \mathbb{Z}[H]^{\oplus 9}\) by
\(s(h-1) = (\dots, 0, \underset{\text{coord } h}{-1}, 0, \dots)\) for
\(h \neq 1\), extended \(\mathbb{Z}\)-linearly.
Then \(\pi(s(h-1)) = -1 - \varepsilon(-1)h = h-1\).

\medskip\noindent
\textbf{Step 3: the \(2\)-cocycle.}
Set \(\widetilde{\psi}(h) = s(h-1)\).  For \(h \neq 1\) and \(gh \neq g\):
\begin{align*}
  \widetilde{\psi}(g)   &\text{ has entry } -1 \text{ in coordinate } g,\\
  \widetilde{\psi}(h)   &\text{ has entry } -1 \text{ in coordinate } h,\\
g \cdot \widetilde{\psi}(h) &\text{ has entry } -g \text{ in coordinate } gh,\\
  \widetilde{\psi}(gh)  &\text{ has entry } -1 \text{ in coordinate } gh.
\end{align*}
Hence for \(h \neq 1\):
\begin{equation}
\boxed{\;
z(g,h) = \bigl(\; \underset{\text{coord }g}{-1},\;
\underset{\text{coord }gh}{1-g},\; 0,\dots,0 \;\bigr) \;\in\; M|_H
\;}.
\end{equation}
For \(h = 1\), \(z(g,1) = 0\).  One verifies \(\pi(z(g,h)) = 0\) directly.
The class \([z]\) generates \(C_9\) since \([\psi]\) generates
\(\coh^1(H, I[H])\) and \(\partial\) is an isomorphism.  The formula depends
on the choice of basis in \(\mathbb{Z}[H]^{\oplus 9}\) and on the splitting
\(s\); different choices yield cohomologous cocycles.

%%%%%%%%%%%%%%%%%%%%%%%%%%%%%%%%%%%%%%%%%%%%%%%%%%%%%%%%%%%%%%%%%%%%%%%%%%%
% [2026-10-01 用户指令] 附录 A 子节标题降级为段标题（\subsection → \paragraph；
% \ref{sec:gap} 仅在被注释行内出现，无 live 引用）
\paragraph{Computational Verification}
\label{sec:gap}
%%%%%%%%%%%%%%%%%%%%%%%%%%%%%%%%%%%%%%%%%%%%%%%%%%%%%%%%%%%%%%%%%%%%%%%%%%%

The cohomological readings used above were checked by machine.  The two
programs are collected in the accompanying code package, described in the
Data availability statement, and are not reproduced here.

\medskip\noindent
The two programs are independent of each other: one is written in Python
using only the standard library, the other in \textsf{GAP} using the
\textsf{HAP} package.  They use different resolutions of
\(\mathbb{Z}\) over \(\mathbb{Z}[H]\) --- a K\"unneth resolution assembled
from the periodic resolutions of the two cyclic factors, and the
inhomogeneous cochain complex --- and independent Smith normal form
implementations.  The Python program also computes
\(\coh^2(H, M|_H)\) directly on the rank-\(73\) lattice \(M|_H\), without
passing through the long exact sequence, and it prints for every reading
unimodular matrices \(U, V\) and a diagonal \(D\) with \(UAV = D\), so
that the computation can be checked with any computer algebra system.  It
reports no value unless a built-in self-test, with positive controls
(independently known values) and negative controls, passes.

\medskip\noindent
For \(\mathbb{Z}[H]\) the ranks of the differentials \((8, 73)\) give
\(\dim(\ker d^1) - \operatorname{rank}(d^0) = 8-8 = 0\), confirming
\(\coh^1(H,\mathbb{Z}[H])=0\).  For \(I[H]\), the ranks \((8,64)\) give free
rank~0, and the augmentation exact sequence yields the torsion subgroup
\(\coh^1(H,I[H]) \cong \mathbb{Z}/9\mathbb{Z} \cong C_9\).  The GAP
computation gives \(\coh^2(H,\mathbb{Z}) \cong C_3 \times C_3\).  By
additivity \cite[Proposition~3.3.4]{Weibel1994},
\(\coh^i(H,\mathbb{Z}[H]^{\oplus 9}) = 0\) for \(i=1,2\).  The Procesi
short exact sequence
\[
0 \longrightarrow M|_H \longrightarrow \mathbb{Z}[H]^{\oplus 9}
\longrightarrow I[H] \longrightarrow 0
\]
then forces the connecting homomorphism
\(\partial : \coh^1(H,I[H]) \to \coh^2(H,M|_H)\) to be an isomorphism,
giving \(\coh^2(H,M|_H) \cong C_9\).  The exponent~\(9\) is incompatible
with stable permutation (all permutation lattices over \(C_3\times C_3\)
have \(\coh^2\)-exponent dividing~\(3\) by
Proposition~\ref{prop:perm-exponent}); hence \(M|_H\) is not stably permutation,
and by Theorem~\ref{thm:endo-miyata}, \(Z_H(F,9)\) is not stably rational (hence
not rational) over \(F\) in characteristic~\(0\).

\medskip\noindent
The same programs with \(p = 5\) give \(|H| = 25\) and
\(\operatorname{rank}(d^0_{\mathbb{Z}[H]}) = 24\),
\(\operatorname{rank}(d^1_{\mathbb{Z}[H]}) = 601\),
\(\operatorname{rank}(d^0_{I[H]}) = 24\),
\(\operatorname{rank}(d^1_{I[H]}) = 576\) with torsion
\(\mathbb{Z}/25\mathbb{Z}\), and \(\coh^2(H,\mathbb{Z}) = [5,5]\).  Hence
\(\coh^2(H,M|_H) \cong C_{25}\), of exponent \(25 > 5\); the exponent \(5\)
is the largest occurring for permutation lattices in this case.

\noindent The pattern
\[
\operatorname{rank}(d^0_{\mathbb{Z}[H]})=p^2-1,\qquad
\operatorname{rank}(d^1_{\mathbb{Z}[H]})=p^4-p^2+1,\qquad
\operatorname{rank}(d^1_{I[H]})=(p^2-1)^2
\]
holds for both \(p=3\) and \(p=5\), with
\(\coh^1(H,I[H]) \cong C_{p^2}\) in both cases; the same
connecting-homomorphism argument as above then gives
\(\coh^2(H,M|_H)\cong C_{p^2}\), of exponent~\(p^2\).  For general odd~\(p\)
the formal cohomological ingredients are identical (see the proof of
Theorem~\ref{thm:Cp-general}); the only computational bottleneck is the
Smith normal form of \(d^1_{\mathbb{Z}[H]}\), of size \(p^6 \times p^3\).

\medskip\noindent
\textbf{Computational scaling.}
The Smith normal form of dense integer matrices is the dominant cost.
For the bar resolution, the matrices are:
\begin{itemize}
  \item \(d^0_{\mathbb{Z}[H]}\): size \(p^3 \times p^2\) (rows $\times$ cols)
  \item \(d^1_{\mathbb{Z}[H]}\): size \(p^6 \times p^3\)
  \item \(d^0_{I[H]}\): size \((p^2-1)p \times (p^2-1)\)
  \item \(d^1_{I[H]}\): size \((p^2-1)p^2 \times (p^2-1)p\)
\end{itemize}
The largest matrix has $p^6 \times p^3$ entries.  For $p = 3$ this is
$729 \times 81$ ($59\,049$ entries), for $p = 5$ it is $15\,625 \times 625$
($9.8$ million entries), and for $p = 7$ it would be $117\,649 \times 2\,401$
($282$ million entries).

% With GAP's \texttt{SmithNormalFormIntegerMat}
% (which computes the full SNF over $\mathbb{Z}$, not modulo primes),
% we observe empirical scaling roughly $\sim O(p^9)$ in time and
% $\sim O(p^6)$ in memory.  On our hardware this places $p = 7$ at the
% boundary of feasibility ($\sim$ hours, $\sim$ 30--50\,GB), and $p = 11$
% beyond reach without modular or sparse techniques.  Since the formal
% proof of Theorem~\ref{thm:Cp-general} covers all odd~$p$, computational
% verification for larger~$p$ is not required for mathematical
% correctness; it serves only as a sanity check on the pattern.

% \medskip\noindent \textbf{Roadmap for
% general~\(p\).}\label{sec:generalization-roadmap} For an odd prime~\(p\), the
% following steps are required to complete the generalization: \begin{enumerate}
% \item[(i)] Construct the Procesi kernel \(M|_H\) as a \(\mathbb{Z}\)-free
% \(\mathbb{Z}[H]\)-lattice via the exact sequence \[ 0 \longrightarrow M|_H
% \longrightarrow \mathbb{Z}[H]^{\oplus p^2} \longrightarrow I[H]
% \longrightarrow 0 ; \] \(\mathbb{Z}\)-freeness is automatic (submodule of a
% free \(\mathbb{Z}\)-module) \cite[Ch.VII, \S3, Th.1, Cor.2]{Bourbaki2007b}.
% \item[(ii)] Verify \(\coh^1(H, I[H]) \cong C_{p^2}\) (augmentation sequence)
% and \(\coh^1(H, \mathbb{Z}[H]) = 0\) (Shapiro's lemma). \item[(iii)] The
% connecting homomorphism \(\partial : \coh^1(H, I[H]) \to \coh^2(H, M|_H)\) is
% an integral isomorphism because \(\coh^i(H, \mathbb{Z}[H]^{\oplus p^2})=0\)
% for \(i=1,2\) (Shapiro + additivity); no class-group or \(p\)-adic integrality
% issues arise. \item[(iv)] Hence \(\coh^2(H, M|_H) \cong C_{p^2}\) with
% exponent~\(p^2\). For any permutation lattice~\(P\), \(\exp(\coh^2(H, P)) \mid
% p\); since \(p^2 \nmid p\), the lattice \(M|_H\) is not stably permutation.
% \end{enumerate} Steps (i)--(iv) are purely formal and hold for every odd
% prime~\(p\). Computational verification for \(p=3\) (Appendix~\ref{sec:gap})
% and \(p=5\) (Appendix~\ref{sec:gap-p5}) confirms the pattern. A complete proof
% for all odd~\(p\) therefore reduces to verifying that the Procesi sequence
% restricts integrally to \(H = C_p \times C_p \subset S_{p^2}\) --- a statement
% we have proved in \S\ref{sec:procesi} without any dependence on~\(p\).




%%%%%%%%%%%%%%%%%%%%%%%%%%%%%%%%%%%%%%%%%%%%%%%%%%%%%%%%%%%%%%%%%%%%%%%%%%%
%%%%%%%%%%%%%%%%%%%%%%%%%%%%%%%%%%%%%%%%%%%%%%%%%%%%%%%%%%%%%%%%%%%%%%%%%%%

\section{Appendix B: Comparison via the unramified Brauer group}
\label{app:brauer}
%%%%%%%%%%%%%%%%%%%%%%%%%%%%%%%%%%%%%%%%%%%%%%%%%%%%%%%%%%%%%%%%%%%%%%%%%%%

For a multiplicative invariant field \(K = F(M)^G\), the unramified
Brauer group \(\Br_{\mathrm{nr}}(K/F)\) is a stable birational invariant
\cite[Theorem~3.1]{Bogomolov1987b} and the principal known obstruction
to rationality and stable rationality.  Below we examine two
complementary ways to access \(\Br_{\mathrm{nr}}\).

For the Noether problem --- whether \(F(H)^H\) is rational over \(F\),
where \(F(H) = F(x_h : h \in H)\) is the function field of the regular
representation \(\mathbb{Z}[H]\) --- Bogomolov's theorem
\cite[Theorem~3.1]{Bogomolov1987b} computes the unramified Brauer
group as
\[
\Br_{\mathrm{nr}}(K/F) \;\cong\;
\ker\!\Bigl( \coh^2(H, \mathbb{Q}/\mathbb{Z})
\longrightarrow \bigoplus_{A < H} \coh^2(A, \mathbb{Q}/\mathbb{Z}) \Bigr),
\]
where \(A\) runs over all abelian subgroups of \(H\) and the map
is the product of restriction homomorphisms
\cite[\S4, Definition~1]{Bogomolov1987b}.

\begin{theorem}\label{thm:UCT}
For a finite group \(G\) and a trivial \(G\)-module \(M\), there is a
natural short exact sequence
\[
0 \longrightarrow \Ext^1_{\mathbb{Z}}(\mathrm{H}_{n-1}(G, \mathbb{Z}), M)
\longrightarrow \coh^n(G, M)
\longrightarrow \Hom(\mathrm{H}_n(G, \mathbb{Z}), M) \longrightarrow 0
\]
for all \(n \ge 1\) \cite[3.6.5]{Weibel1994}.
If \(M\) is divisible (e.g.\ \(M = \mathbb{Q}/\mathbb{Z}\)), the
sequence splits, giving a direct sum decomposition.
\end{theorem}

\begin{proposition}\label{prop:B0}
\(\Br_{\mathrm{nr}}(F(H)^H/F) \cong C_3\).
\end{proposition}

\begin{proof}
Apply the Universal Coefficient Theorem~\ref{thm:UCT} with
\(M = \mathbb{Q}/\mathbb{Z}\) (divisible, so the sequence splits).
For \(H = C_3 \times C_3\):
\(\mathrm{H}_1(H,\mathbb{Z}) = H \cong C_3 \times C_3\) and
\(\mathrm{H}_2(H,\mathbb{Z}) \cong \bigwedge^2_{\mathbb{Z}} H \cong C_3\)
\cite[Ch.\ V, Theorem~6.4(iii)]{Brown1982}.  Hence
\[
\coh^2(H, \mathbb{Q}/\mathbb{Z}) \cong
\Hom(C_3, \mathbb{Q}/\mathbb{Z}) \oplus
\Ext^1_{\mathbb{Z}}(C_3 \times C_3, \mathbb{Q}/\mathbb{Z})
\cong (C_3)^3 .
\]
For a proper subgroup \(K \cong C_3\):
\(\mathrm{H}_1(K, \mathbb{Z}) \cong C_3\) and \(\mathrm{H}_2(K, \mathbb{Z}) =
0\) (cyclic). By UCT, \(\coh^2(K, \mathbb{Q}/\mathbb{Z}) \cong
\Ext^1_{\mathbb{Z}}(C_3, \mathbb{Q}/\mathbb{Z}) \cong C_3\).

Identify \(\Ext^1_{\mathbb{Z}}(C_3 \times C_3, \mathbb{Q}/\mathbb{Z}) \cong
\Hom(C_3, \mathbb{Q}/\mathbb{Z})^{\oplus 2}\) via the two projections.
For \(K = \langle a^i b^j \rangle\), the restriction on the \(\Ext\) summand
sends \((\varphi_1, \varphi_2) \mapsto i\varphi_1 + j\varphi_2\).  Over the
four cyclic subgroups \(\langle a\rangle, \langle b\rangle,
\langle ab\rangle, \langle ab^2\rangle\) this gives the \(\mathbb{F}_3\)-linear
map
\[
(\varphi_1, \varphi_2) \longmapsto
(\varphi_1,\; \varphi_2,\; \varphi_1+\varphi_2,\; \varphi_1+2\varphi_2)
\in (C_3)^{\oplus 4},
\]
which has rank~\(2\), hence is injective.  The \(\Hom\) summand lies in the
kernel of every restriction because \(\mathrm{H}_2(K, \mathbb{Z}) = 0\).
Therefore the kernel is \(\Hom(\mathrm{H}_2(H,\mathbb{Z}),
\mathbb{Q}/\mathbb{Z})
\cong C_3\).
(The same computation gives \(\Br_{\mathrm{nr}}(F(H)^H/F) \cong C_p\)
for all odd primes~\(p\); see \cite[\S6, Lemma~1]{Bogomolov1987b}.)
\end{proof}

\begin{remark}
Thus \(\Br_{\mathrm{nr}}(F(H)^H/F) \cong C_3\) is non-trivial.
Noether's problem asks whether
\(F(H)^H = F(x_h : h \in H)^H\)
(the multiplicative invariant field of the regular representation
\(\mathbb{Z}[H]\)) is rational over \(F\);
by \cite[Cor.\ of Lem.~2.1]{Bogomolov1987b},
\(F(H)^H\) is therefore not even retract rational.
In contrast, \(Z_H(F,9) = F(M|_H)^H\) is retract rational
(Theorem~\ref{thm:retract}) but not stably rational
(Theorem~\ref{thm:not-stably-rational}).
Thus \(Z_H(F,9)\) and the Noether problem exhibit opposite rationality
behavior, despite both being multiplicative invariant fields of the
same group~\(H\); the integral \(\coh^2\) obstruction, not
\(\Br_{\mathrm{nr}}\), governs \(Z_H(F,9)\).
\end{remark}

\medskip\noindent
\textbf{Hoshi--Kang--Yamasaki.}
For general multiplicative invariant fields, Hoshi, Kang, and
Yamasaki \cite{HKY} have systematically computed unramified Brauer
groups for small groups; their tables cover groups with \(\le 6\)
variables.  For groups of order \(9\), all unramified Brauer obstructions
originate from \(\Br_{\mathrm{nr}} \cong C_3\).  Our field \(Z_H(F,9)\) has
transcendence degree \(73\) and lies outside their range.
Both approaches operate via
\(\Br_{\mathrm{nr}}\), a geometric invariant of the function field;
by contrast, our \(\coh^2\) obstruction is a purely integral
cohomological invariant of the lattice~\(M|_H\).
Theorem~\ref{thm:not-stably-rational} thus exhibits a
multiplicative invariant field whose non-stable-rationality is detected by
an integral cohomological obstruction, which appears to be new in this
context.

\medskip\noindent
% \textbf{Quasi-permutation modules.}
% A \(\mathbb{Z}[H]\)-lattice~\(L\) is \emph{quasi-permutation} if
% \(L \oplus P \cong Q\) for permutation lattices \(P, Q\).
% Endo--Miyata \cite[Theorem~1.2]{EndoMiyata1973b} proved that
% \(F(L)^H\) is stably rational over \(F\) if and only if
% \(L\) is quasi-permutation.  Establishing quasi-permutation
% constructively requires exhibiting the lattices \(P, Q\).
% By contrast, our \(\coh^2\)-exponent obstruction
% (Theorem~\ref{thm:not-stably-perm}) rules out stable permutation
% from a single cohomological computation
% (\(\coh^2(H, M|_H) \cong C_9\), Lemma~\ref{lem:H2-MH}),
% without classifying the lattice.

% \medskip\noindent \textbf{Saltman invariants.} Saltman \cite[\S3]{Saltman1984}
% introduced the invariant \(\Br(F)/\Br_0(F)\) and used it to construct the
% first non-rational stably rational fields. For \(H = C_3 \times C_3\),
% Saltman's retract rationality theorem \cite[Corollary~3.13]{Saltman1984} gives
% retract rationality unconditionally; our \(\coh^2\) obstruction then shows
% that retract rationality is optimal: neither rationality nor stable
% rationality holds.

%%%%%%%%%%%%%%%%%%%%%%%%%%%%%%%%%%%%%%%%%%%%%%%%%%%%%%%%%%%%%%%%%%%%%%%%%%%
%%%%%%%%%%%%%%%%%%%%%%%%%%%%%%%%%%%%%%%%%%%%%%%%%%%%%%%%%%%%%%%%%%%%%%%%%%%

\section{Appendix C: Scope and further directions}
\label{app:scope}
%%%%%%%%%%%%%%%%%%%%%%%%%%%%%%%%%%%%%%%%%%%%%%%%%%%%%%%%%%%%%%%%%%%%%%%%%%%

\medskip\noindent
\textbf{Non-regular embeddings.}
Our proof uses the regular embedding \(H \hookrightarrow S_{p^2}\)
(Proposition~\ref{prop:embedding-independence}) in three
essential ways.  We itemize them to clarify what would change for other
embeddings.

\begin{enumerate}
  \item \emph{Identification \(U|_H \cong \mathbb{Z}[H]\).}
    This relies on the action of \(H\) on \(\Omega = S_{p^2}/S_{p^2-1}\)
    being regular (free and transitive, hence \(H\)-equivariantly isomorphic
    to the left regular action), and is independent of the
    basepoint up to a \(\mathbb{Z}[H]\)-module automorphism
(Proposition~\ref{prop:U-restrict-canonical}). If \(H\) acts on \(\Omega\) with
multiple orbits or with non-trivial stabilizers, then \(U|_H \cong \bigoplus_{K
\le H} m_K \cdot \mathbb{Z}[H/K]\) where the multiplicities \(m_K\) depend on
the orbit structure. By Shapiro's lemma \cite[Theorem~6.3.2]{Weibel1994},
\(\mathbb{Z}[H/K] \cong \Ind_K^H \mathbb{Z}\); the tensor product \(\Ind_K^H
\mathbb{Z} \otimes \Ind_L^H \mathbb{Z}\) decomposes via Mackey's formula
\cite[(10.18)]{CurtisReiner1981} as a direct sum of \(\Ind_{K \cap gLg^{-1}}^H
\mathbb{Z}\) over double cosets. Since \(H\) is abelian, \(gLg^{-1}=L\) and
every intersection is \(K \cap L\). Applying Shapiro's lemma again yields
    \[
    \coh^i(H, U|_H \otimes U|_H) \cong \bigoplus_{K,L} m_K m_L \,
    \coh^i(K \cap L, \mathbb{Z}),
    \]
    generally non-zero for \(i>0\), so the
    connecting-homomorphism isomorphism
    \(\coh^2(H, M|_H) \cong \coh^1(H, I[H])\) would
    no longer be an isomorphism, so
    \(\coh^2(H, M|_H) \cong C_9\) (Lemma~\ref{lem:H2-MH}) would fail.

  \item \emph{Breakdown of \(\mathbb{Z}[H] \otimes \mathbb{Z}[H]
    \cong \mathbb{Z}[H]^{\oplus p^2}\).}
    In the regular case
    \(\mathbb{Z}[H] \otimes_{\mathbb{Z}} \mathbb{Z}[H] \cong
    \mathbb{Z}[H]^{\oplus p^2}\)
    (Proposition~\ref{prop:tensor-reg-ZH}), which relies on
    \(\Res_1^H \Ind_1^H \mathbb{Z} \cong \mathbb{Z}^{\oplus p^2}\).
    For a non-regular \(\mathbb{Z}[H/K]\), however, the tensor product
    \(\mathbb{Z}[H/K] \otimes \mathbb{Z}[H/L]\) decomposes via
    Mackey's formula \cite[(10.18)]{CurtisReiner1981} into summands
    indexed by double cosets; the simple rank \(p^2\) is replaced by
    a sum of \([H:K \cap L]\) terms, so the cohomological analysis
    changes.

  \item \emph{Surjectivity \(\mathbb{Z}[H]^{\oplus p^2} \to I[H]\).}
    This is the rightmost map in the Procesi sequence
    (Proposition~\ref{prop:procesi-H}), where \(I[H]\) is generated
    as a \(\mathbb{Z}[H]\)-module by \(h-1\) and the surjection from
    \(p^2\) copies of \(\mathbb{Z}[H]\) is guaranteed by the regular
    action.  For a non-regular permutation lattice, the image of the
    Procesi map \(\varphi\) is a submodule of \(I[\Omega]\) that may
    be strictly smaller; the exact sequence of
    Proposition~\ref{prop:procesi-H} would no longer terminate in
    \(I[H]\), requiring analysis of the specific orbit decomposition.
\end{enumerate}

In summary, for a non-regular embedding the short exact sequence
\[
0 \longrightarrow M|_H \longrightarrow P_1 \longrightarrow I[\Omega]
\longrightarrow 0
\]
(with \(P_1\) a permutation lattice not necessarily of the form
\(\mathbb{Z}[H]^{\oplus n}\)) remains exact, but the cohomological
analysis of \(P_1\) becomes
substantially more involved and the simple exponent obstruction may
no longer yield a definitive answer.  Computing \(\coh^2(H, M|_H)\) for
one or two non-regular embeddings of \(C_3 \times C_3\) into
\(S_n\) (e.g.\ the intransitive embedding \(C_3 \times C_3 \subset
S_3 \times S_3 \subset S_6\)) is an \textbf{open} problem.

\medskip\noindent
\textbf{Non-abelian \(p\)-groups.}
For non-abelian groups the situation is more complex.  The subgroup
cohomology classification used in Proposition~\ref{prop:perm-exponent}
extends: if \(G\) is any finite group and \(P = \bigoplus_i
\mathbb{Z}[G/K_i]\) is a permutation lattice, then
\(\exp(\coh^2(G, P)) \mid \max_{K \le G} \exp(\coh^2(K, \mathbb{Z}))\) by
Shapiro's lemma.  The obstruction therefore applies whenever one can
find a natural \(G\)-lattice~\(M\) with
\(\exp(\coh^2(G, M)) > \max_{K \le G} \exp(\coh^2(K, \mathbb{Z}))\).

For the Heisenberg group
\(G = \langle x,y \mid x^p = y^p = 1,\ [x,y] \text{ central},\
[x,y]^p = 1\rangle\), with \(p\) odd.  Write \(c = [x,y]\).  The
quotient \(G/\langle c\rangle\) is abelian and generated by two elements
of order dividing \(p\), so \(|G| \le p^3\); and \(G\) maps onto the
group of upper unitriangular \(3 \times 3\) matrices over
\(\mathbb{F}_p\), so \(|G| \ge p^3\).  Thus \(|G| = p^3\), with centre and
derived subgroup both \(\langle c\rangle\).  Since \(G\) has nilpotency
class \(2\) and \(p\) is odd, the power map \(g \mapsto g^p\) is a
homomorphism; it vanishes on \(x\) and \(y\), hence on \(G\), so
\(\exp G = p\).  Therefore \(G\) is the extraspecial \(p\)-group of order
\(p^3\) and exponent \(p\), and
\(G^{\mathrm{ab}} \cong C_p \times C_p\) and
\(\mathrm{H}_2(G, \mathbb{Z}) \cong C_p \times C_p\)
\cite[Corollary~IV.4.16(b)]{BeylTappe1982}.  By the Universal Coefficient
Theorem~\ref{thm:UCT},
\(\coh^2(G, \mathbb{Z}) \cong \Ext^1_{\mathbb{Z}}(G^{\mathrm{ab}}, \mathbb{Z})
\cong C_p \times C_p\).  Hence the maximal permutation exponent is
\(p\).  Whether a natural~\(M\) with \(\exp(\coh^2(G, M)) = p^2\) exists
for such groups (e.g.\ from a Procesi-like construction or another natural
\(G\)-lattice) is an \textbf{open} question.

% [2026-10-01 第二十一改，用户指令] flasque 整块自 §Preliminaries 移入本附录
% （原址留注释存档；全篇唯一实质使用者即下方 Br 段）。以下为原样文字：
\begin{definition}\label{def:flasque}
A \(\mathbb{Z}[G]\)-lattice \(F\) is \emph{flasque} if
\(\hat{\coh}^{-1}(K, F) = 0\) for all subgroups \(K \le G\)
(equivalently, \(\Ext^1_G(F, P) = 0\) for every permutation lattice~\(P\))
\cite[\S1, Lemme~9]{CTS77}. Here \(\hat{\coh}^{-1}\) denotes
Tate cohomology in degree \(-1\); for a finite group~\(G\) and a
\(\mathbb{Z}[G]\)-lattice~\(M\),
\(\hat{\coh}^{-1}(G, M) = \ker(N_G) / I[G]M\), where
\(N_G = \sum_{g \in G} g\) is the norm element and
\(I[G] = \ker(\mathbb{Z}[G] \twoheadrightarrow \mathbb{Z})\) is the
augmentation ideal.
\end{definition}

\begin{definition}\label{def:flasque-class-group}
Two \(\mathbb{Z}[G]\)-lattices \(M, N\) are called \emph{similar}
(written \(M \sim N\)) if \(M \oplus S_1 \cong N \oplus S_2\) for
some permutation lattices \(S_1, S_2\) (Definition~\ref{def:lattice}).
This is an equivalence relation; denote the similarity class of \(M\)
by \([M]\). Let
\[
C(G)/S(G) = \{ [M] : M \text{ a finitely generated }
\mathbb{Z}[G]\text{-lattice} \},
\]
a commutative monoid under direct sum, with zero element the class
of permutation lattices. Its maximal subgroup is
\[
D(G)/S(G) = \{ [M] \in C(G)/S(G) : M \text{ is stably permutation}\},
\]
which is the semigroup \(C(G)/S(G)\) of \cite[\S4.7]{Vos98}.
\end{definition}

\begin{proposition}\label{prop:flasque-resolution}
Every \(\mathbb{Z}[G]\)-lattice~\(M\) has a \emph{flasque resolution}
\[
0 \longrightarrow M \longrightarrow P \longrightarrow F \longrightarrow 0
\]
with \(P\) permutation (Definition~\ref{def:lattice}) and \(F\) flasque
(Definition~\ref{def:flasque})
\cite[\S1, Lemme~3]{CTS77}. The similarity class
\([F] \in D(G)/S(G)\) (Definition~\ref{def:flasque-class-group})
depends only on \(M\)
\cite[\S1, Lemme~5]{CTS77}; we denote it by
\(\rho(M) = [F]\) and call it the
\emph{flasque class} of \(M\).
By construction, \(M\) is stably permutation iff \(\rho(M) = 0\)
\cite[\S1, Lemmes~7 and~8]{CTS77}; this is the defining property
of the group \(D(G)/S(G)\).
\end{proposition}

\medskip\noindent
\textbf{Unramified Brauer group.}
Let
\[
0 \longrightarrow M|_H \longrightarrow P \longrightarrow F \longrightarrow 0
\]
be a flasque resolution of \(M|_H\)
(Proposition~\ref{prop:flasque-resolution}).
From the long exact sequence:
\[
\coh^1(H, P) \longrightarrow \coh^1(H, F) \longrightarrow \coh^2(H, M|_H)
\longrightarrow \coh^2(H, P),
\]
where \(\coh^1(H, P)=0\) (Shapiro's lemma \cite[Theorem~6.3.2]{Weibel1994}),
\(\coh^2(H, M|_H) \cong C_9\) (Lemma~\ref{lem:H2-MH}), and
\(\exp(\coh^2(H, P)) \mid 3\) (Proposition~\ref{prop:perm-exponent}).  Hence
\(\ker(\coh^2(M|_H) \to \coh^2(P)) \cong \coh^1(H, F)\) contains the unique
subgroup \(C_3 \subset C_9\); the containment can be \(C_3\) or \(C_9\).

The field \(Z_H(F,9) \cong F(M|_H)^H\) is the function field of the
\(H\)-torus with character group \(M|_H\).  Let \(X\) be a smooth
\(F\)-compactification of this torus.  By
\cite[Proposition~9.5(ii)]{CTS87},
\[
\Br(X)/\Br(F) = \coh^1(H, \rho(M|_H)) = \coh^1(H, F),
\]
where \(\rho(M|_H) = [F] \in D(H)/S(H)\) is the flasque class and
the second equality follows from the proof of
\cite[Proposition~9.5]{CTS87} via the flasque resolution above.
Since \(\Br_{\mathrm{nr}}(Z_H(F,9)) = \Br(X)\) for a smooth
compactification~\cite[\S9]{CTS87},
\[
\Br_{\mathrm{nr}}(Z_H(F,9)) / \Br(F) \;\cong\; \coh^1(H, F)
\;\supseteq\; C_3 .
\]
Thus \(\Br_{\mathrm{nr}}(Z_H(F,9)) \neq \Br(F)\); the unramified Brauer group
contains at least a \(C_3\) summand beyond the base field.  This is the
portion of the integral \(\coh^2\) obstruction \(C_9\) that is detected by a
classical stable birational invariant.  The exact value (\(C_3\) or \(C_9\))
remains \textbf{open} pending an explicit flasque resolution.

\medskip\noindent
\textbf{Characteristic-\(0\) dependence.}
The cohomological computations in \S\ref{sec:cohom} (the isomorphisms
\(\coh^2(H, M|_H) \cong C_9\) and \(\exp(\coh^2(H, P)) \mid 3\)) are purely
integral and valid whenever \(\Char F \neq 3\).  The geometric
conclusion that \(Z_H(F,9)\) is not stably rational
(Theorem~\ref{thm:not-stably-rational}) additionally requires
characteristic~\(0\) (the criterion is stated in characteristic \(0\)):
Theorem~\ref{thm:endo-miyata}, used as a hypothesis, together with the
characteristic-\(0\) Procesi isomorphism.  Whether
the non-stable-rationality conclusion extends to positive characteristic
(not dividing~\(9\)) remains an \textbf{open} question.
Theorem~\ref{thm:retract} (retract rationality) holds over any field, because
\(H = C_3 \times C_3\) is abelian of odd order.

%%%%%%%%%%%%%%%%%%%%%%%%%%%%%%%%%%%%%%%%%%%%%%%%%%%%%%%%%%%%%%%%%%%%%%%%%%%
%%%%%%%%%%%%%%%%%%%%%%%%%%%%%%%%%%%%%%%%%%%%%%%%%%%%%%%%%%%%%%%%%%%%%%%%%%%


\begin{filecontents}{wu.bib}
@article{ABGV2011,
	author = {Auel, A. and Brussel, E. and Garibaldi, S. and Vishne, U.},
	title = {Open problems on central simple algebras.},
	fjournal = {Transformation Groups},
	journal = {Transform. Groups},
	issn = {1083-4362},
	volume = {16},
	number = {1},
	pages = {219--264},
	year = {2011},
	language = {English},
	doi = {10.1007/s00031-011-9119-8},
	keywords = {16K20,16K50,16-02,00A07},
	zbMATH = {5898840},
	Zbl = {1230.16016}
}	
	
	

@misc{CurtisReiner1981,
	author = {Curtis, C. W. and Reiner, I.},
title = {Methods of representation theory, with applications to finite groups
and orders. {Vol}. {I}}, year = {1981}, language = {English}, howpublished =
{Pure and {Applied} {Mathematics}. {A} {Wiley}-{Interscience} {Publication}.
{New} {York} etc.: {John} {Wiley} \& {Sons}. {XXI}, 819 p. {{\textsterling}}
40.70 (1981).}, keywords =
{20-02,16-02,20Cxx,16H05,16Gxx,16U30,20C05,20C10,20C15,20C20,20C25,16P10},
zbMATH = {3736029}, Zbl = {0469.20001} }

@book{Brown1982,
    author = {Brown, K. S.},
    title = {Cohomology of groups},
    fseries = {Graduate Texts in Mathematics},
    series = {Grad. Texts Math.},
    issn = {0072-5285},
    volume = {87},
    year = {1982},
    publisher = {Springer, Cham},
    language = {English},
    keywords = {20J05,20-02,18-02,55-02,20J06},
    zbMATH = {3935317},
    Zbl = {0584.20036}
}

@book{Weibel1994,
    author = {Weibel, C. A.},
    title = {An introduction to homological algebra},
    fseries = {Cambridge Studies in Advanced Mathematics},
    series = {Camb. Stud. Adv. Math.},
    volume = {38},
    isbn = {0-521-43500-5},
    year = {1994},
    publisher = {Cambridge: Cambridge University Press},
    language = {English},
    keywords = {18-01,18Gxx,16Exx,18E30,20F99,17B55,55Uxx},
    zbMATH = {595200},
    Zbl = {0797.18001}
}



@book{HKY,
    author = {Hoshi, A. and Kang, M. and Yamasaki, A.},
    title = {Multiplicative invariant fields of dimension {{\(\le 6\)}}},
    fseries = {Memoirs of the American Mathematical Society},
    series = {Mem. Am. Math. Soc.},
    issn = {0065-9266},
    volume = {1403},
    isbn = {978-1-4704-6022-8; 978-1-4704-7404-1},
    year = {2023},
    publisher = {Providence, RI: American Mathematical Society (AMS)},
    language = {English},
    doi = {10.1090/memo/1403},
    keywords = {14-02,14E08,20C10,20J06,14F22},
    zbMATH = {7653284},
    Zbl = {1530.14001}
}

@book{CurtisReiner1987,
    author = {Curtis, C. W. and Reiner, I.},
title = {Methods of representation theory. {With} applications to finite groups
and orders. {Vol}. {II}}, year = {1987}, language = {English}, howpublished =
{Pure and Applied Mathematics. A Wiley-Interscience Publication. New York etc.:
John Wiley \& Sons. XIV, 951 p.}, keywords =
{20-02,16-02,20Cxx,16H05,16Gxx,16U30,20C05,20C10,20C15,20C20,20C25,16P10},
zbMATH = {3736029}, Zbl = {0469.20001} }

@article{CTS77,
	author = {Colliot-Th{\'e}l{\`e}ne, J.-L. and Sansuc, J.-J.},
	title = {La {R}-{\'e}quivalence sur les tores},
fjournal = {Annales Scientifiques de l'{\'E}cole Normale Sup{\'e}rieure.
Quatri{\`e}me S{\'e}rie}, journal = {Ann. Sci. {\'E}c. Norm. Sup{\'e}r. (4)},
issn = {0012-9593}, volume = {10}, pages = {175--229}, year = {1977}, language =
{French}, doi = {10.24033/asens.1325}, keywords = {14G05,14F20,14C15}, url =
{https://eudml.org/doc/81994}, zbMATH = {3554404}, Zbl = {0356.14007} }
	
@article{CTS87,
	author = {Colliot-Th{\'e}l{\`e}ne, J.-L. and Sansuc, J.-J.},
	title = {Principal homogeneous spaces under flasque tori; applications},
	fjournal = {Journal of Algebra},
	journal = {J. Algebra},
	issn = {0021-8693},
	volume = {106},
	pages = {148--205},
	year = {1987},
	language = {English},
	doi = {10.1016/0021-8693(87)90026-3},
	keywords = {14F20,14M17,13B02,14L30,20G10},
	zbMATH = {3961811},
	Zbl = {0597.14014}
}

% @article{Bogomolov1987a, author = {Bogomolov, F. A.}, title = {The stable
% rationality of quotient spaces for simply connected groups}, fjournal =
% {Mathematics of the USSR, Sbornik}, journal = {Math. USSR, Sb.}, issn =
% {0025-5734}, volume = {58}, pages = {1--14}, year = {1987}, language =
% {English}, doi = {10.1070/SM1987v058n01ABEH003089}, keywords =
% {14M20,14L30,14M17,14L24}, zbMATH = {4114796}, Zbl = {0681.14030} }

@article{Bogomolov1987b,
	author = {Bogomolov, F. A.},
	title = {Brauer groups of factor spaces of linear representations},
	fjournal = {Izvestiya Akademii Nauk SSSR. Seriya Matematicheskaya},
	journal = {Izv. Akad. Nauk SSSR, Ser. Mat.},
	issn = {0373-2436},
	volume = {51},
	number = {3},
	pages = {485--516},
	year = {1987},
	language = {Russian},
	keywords = {14L24,20G05,14F22,14L30},
	zbMATH = {4043985},
	Zbl = {0641.14005}
}

@book{Bourbaki2006,
	author = {Bourbaki, Nicolas},
title = {{\'E}l{\'e}ments de math{\'e}matique. {Alg{\`e}bre} commutative.
{Chapitres} 1 {\`a} 4.}, edition = {Reprint of the 1985 original}, isbn =
{3-540-33937-X}, year = {2006}, publisher = {Berlin: Springer}, language =
{French}, doi = {10.1007/978-3-540-33976-2}, zbMATH = {5071153}, }

@book{Bourbaki2007a,
	author = {Bourbaki, Nicolas},
title = {{\'E}l{\'e}ments de math{\'e}matique. {Alg{\`e}bre}. {Chapitres} 1
{\`a} 3}, edition = {Reprint of the 1970 original}, isbn = {3-540-33849-7}, year
= {2007}, publisher = {Berlin: Springer}, language = {French}, zbMATH =
{5069335}, }

@book{Bourbaki2007b,
	author = {Bourbaki, Nicolas},
title = {{\'E}l{\'e}ments de math{\'e}matique. {Alg{\`e}bre}. {Chapitres} 4
{\`a} 7}, edition = {Reprint of the 1981 original}, isbn = {3-540-34398-9}, year
= {2007}, publisher = {Berlin: Springer}, language = {French}, doi =
{10.1007/978-3-540-34499-5}, zbMATH = {5069336}, }

@article{EndoMiyata1973a,
	author = {Endo, Shizuo and Miyata, Takehiko},
	title = {Invariants of finite {Abelian} groups},
	fjournal = {Journal of the Mathematical Society of Japan},
	journal = {J. Math. Soc. Japan},
	issn = {0025-5645},
	volume = {25},
	pages = {7--26},
	year = {1973},
	language = {English},
	doi = {10.2969/jmsj/02510007},
	keywords = {20C10,15A72,20K99},
	zbMATH = {3387583},
	Zbl = {0245.20007}
}

@article{EndoMiyata1973b,
	author = {Endo, S. and Miyata, T.},
	title = {Quasi-permutation modules over finite groups},
	fjournal = {Journal of the Mathematical Society of Japan},
	journal = {J. Math. Soc. Japan},
	issn = {0025-5645},
	volume = {25},
	pages = {397--421},
	year = {1973},
	language = {English},
	doi = {10.2969/jmsj/02530397},
	keywords = {20C05,16S34,20C10},
	zbMATH = {3435716},
	Zbl = {0277.20005}
}

% @article{EndoMiyata1974,
% 	author = {Endo, S. and Miyata, T.},
% 	title = {Quasi-permutation modules over finite groups. {II}},
% 	fjournal = {Journal of the Mathematical Society of Japan},
% 	journal = {J. Math. Soc. Japan},
% 	issn = {0025-5645},
% 	volume = {26},
% 	pages = {698--713},
% 	year = {1974},
% 	language = {English},
% 	doi = {10.2969/jmsj/02640698},
% 	keywords = {20C10},
% 	zbMATH = {3448770},
% 	Zbl = {0286.20007}
% }

@article{Swan1960,
    author = {Swan, R. G.},
    title = {Induced representations and projective modules},
    fjournal = {Annals of Mathematics. Second Series},
    journal = {Ann. Math. (2)},
    issn = {0003-486X},
    volume = {71},
    pages = {552--578},
    year = {1960},
    language = {English},
    doi = {10.2307/1969944},
    Zbl = {0104.25102}
}

% @article{Swan1969,
%     author = {Swan, R. G.},
%     title = {Invariant rational functions and a problem of {Steenrod}},
%     fjournal = {Inventiones Mathematicae},
%     journal = {Invent. Math.},
%     issn = {0020-9910},
%     volume = {7},
%     pages = {148--158},
%     year = {1969},
%     language = {English},
%     doi = {10.1007/BF01389798},
%     Zbl = {0174.05203}
% }

@article{Lenstra1974,
    author = {Lenstra, H. W., Jr.},
    title = {Rational functions invariant under a finite abelian group},
    fjournal = {Inventiones Mathematicae},
    journal = {Invent. Math.},
    issn = {0020-9910},
    volume = {25},
    pages = {299--325},
    year = {1974},
    language = {English},
    doi = {10.1007/BF01389732},
    Zbl = {0292.12010}
}

@book{NSW2020,
    author = {Neukirch, J{\"u}rgen and Schmidt, Alexander and Wingberg, Kay},
    title = {Cohomology of number fields},
    edition = {2nd ed.},
    fseries = {Grundlehren der Mathematischen Wissenschaften},
    series = {Grundlehren Math. Wiss.},
    issn = {0072-7830},
    volume = {323},
    isbn = {978-3-540-37888-4},
    year = {2008},
    publisher = {Berlin: Springer},
    language = {English},
    keywords = {11-02,11R34,11R23,11S25,18G10,18G20,20J05,11R37},
    zbMATH = {5162349},
    Zbl = {1136.11001}
}

% @article{Hoshi2020, author = {Hoshi, A.}, title = {Noether's problem and
% rationality problem for multiplicative invariant fields: a survey}, fjournal =
% {RIMS K{\^o}ky{\^u}roku Bessatsu}, journal = {RIMS K{\^o}ky{\^u}roku
% Bessatsu}, issn = {1881-6193}, volume = {B77}, pages = {29--53}, year =
% {2020}, language = {English}, keywords =
% {12F12,13A50,14E08,14F22,11E72,14M20,20J06,12-02}, url =
% {repository.kulib.kyoto-u.ac.jp/dspace/bitstream/2433/260609/1/B77-01.pdf},
% zbMATH = {7303982}, Zbl = {1455.12004} }

@article{EndoMiyata1975,
	author = {Endo, S. and Miyata, T.},
	title = {On a classification of the function fields of algebraic tori},
	fjournal = {Nagoya Mathematical Journal},
	journal = {Nagoya Math. J.},
	issn = {0027-7630},
	volume = {56},
	pages = {85--104},
	year = {1975},
	language = {English},
	doi = {10.1017/S0027763000016408},
	keywords = {14H05,20C05,14G99},
	zbMATH = {3470565},
	Zbl = {0301.14008}
}

@article{Formanek2002,
	author = {Formanek, E.},
	title = {The ring of generic matrices.},
	fjournal = {Journal of Algebra},
	journal = {J. Algebra},
	issn = {0021-8693},
	volume = {258},
	number = {1},
	pages = {310--320},
	year = {2002},
	language = {English},
	doi = {10.1016/S0021-8693(02)00506-9},
	keywords = {16R30,16R20,15-02,16-02},
	zbMATH = {1868090},
	Zbl = {1047.16013}
}	

@article{Procesi1967,
	author = {Procesi, C.},
	title = {Non-commutative affine rings},
fjournal = {Atti della Accademia Nazionale dei Lincei. Memorie. Classe di
Scienze Fisiche, Matematiche e Naturali. Serie VIII. Sezione I (Matematica,
Meccanica, Astronomia, Geodesia e Geofisica)}, journal = {Atti Accad. Naz.
Lincei, Mem., Cl. Sci. Fis. Mat. Nat., VIII. Ser., Sez. I}, issn = {0365-0286},
volume = {8}, pages = {239--255}, year = {1967}, language = {English}, zbMATH =
{3323928}, Zbl = {0204.04802} }


@article{Saltman1984,
	author = {Saltman, D. J.},
	title = {Retract rational fields and cyclic {Galois} extensions},
	fjournal = {Israel Journal of Mathematics},
	journal = {Isr. J. Math.},
	issn = {0021-2172},
	volume = {47},
	pages = {165--215},
	year = {1984},
	language = {English},
	doi = {10.1007/BF02760515},
	keywords = {13B05,12F10},
	zbMATH = {3869502},
	Zbl = {0546.14013}
}	

@article{Saltman1987,
	author = {Saltman, D. J.},
	title = {Multiplicative field invariants},
	fjournal = {Journal of Algebra},
	journal = {J. Algebra},
	issn = {0021-8693},
	volume = {106},
	pages = {221--238},
	year = {1987},
	language = {English},
	doi = {10.1016/0021-8693(87)90028-7},
	keywords = {14F22,11R58,14M20,12G05,12F20,14L30,14L24,12F10},
	zbMATH = {4008526},
	Zbl = {0622.13002}
}	

	
@book{Vos98,
	author = {Voskresenski{\u{\i}}, V. E.},
title = {Algebraic groups and their birational invariants. {Transl}. from the
original {Russsian} manuscript by {Boris} {Kunyavskii}.}, edition = {Rev.
version of `{Algebraic} tori', {Nauka} 1977}, fseries = {Translations of
Mathematical Monographs}, series = {Transl. Math. Monogr.}, issn = {0065-9282},
volume = {179}, isbn = {0-8218-0905-9}, year = {1998}, publisher = {Providence,
RI: American Mathematical Society}, language = {English}, keywords =
{14L35,20-02,14-02,14E05,20G30,11E72,14G27,14L40}, zbMATH = {1204077}, Zbl =
{0974.14034} } @book{BeylTappe1982, author = {Beyl, F. Rudolf and Tappe,
J{\"u}rgen}, title = {Group extensions, representations, and the {Schur}
multiplicator}, fseries = {Lecture Notes in Mathematics}, series = {Lect. Notes
Math.}, issn = {0075-8434}, volume = {958}, year = {1982}, publisher =
{Springer, Cham}, language = {English}, doi = {10.1007/bfb0067022}, keywords =
{20-02,20J05,20C25,20E22,20C15,20E10,20F05,20C20}, zbMATH = {3865539}, Zbl =
{0544.20001} }
\end{filecontents}
%\printbibliography
\bibliographystyle{plain}

\bibliography{wu}
\end{document}
